% ==========================================================================
% solucion-espana.tex — Solución de máximos del examen de ejemplo de España
% Guía del dictamen (1D01), entregable 5. Autor: Víctor Gutiérrez Marcos.
% Texto de examen (raya morada a la izquierda, entorno «examen») y cuadros
% complementarios (fondo gris y raya dorada, entorno «complemento»), una
% sección por pregunta (secciones/p01.tex … p10.tex).
% Cifras: cuadernillo del examen (../datos, consulta 26/10/2026); cálculos en
% calculos.py → calculos.txt y datos/; gráficos propios en graficos/s*.tex;
% infografías de la guía en ../../graficos/. Compilación: README.md.
% ==========================================================================
\documentclass[11pt]{article}
\usepackage{tcee}
\addbibresource{../../1D01.bib}
\tcetema[Solución del examen de ejemplo: España]{Primer ejercicio}{Solución de máximos del examen de ejemplo del dictamen: España}
\tcefecha{26/10/2026}
\makeatletter
\renewcommand{\tcee@nombrepags}{Solución}
\makeatother
\graphicspath{{graficos/}{../../graficos/}}

% --- Cajas propias de la solución ------------------------------------------------
% Enunciado de cada pregunta, con su peso, su tiempo y su extensión
\newtcolorbox{enunciado}{enhanced,breakable,colback=tceesalvia,colframe=tceesalviaoscuro,
  boxrule=0.5pt,arc=2pt,fontupper=\small,left=8pt,right=8pt,top=4pt,bottom=4pt}
% Texto de examen: lo que escribiría el opositor (raya morada a la izquierda)
\newtcolorbox{examen}{enhanced,breakable,frame hidden,colback=white,interior hidden,
  borderline west={1.6pt}{0pt}{tceemorado},left=10pt,right=0pt,top=1pt,bottom=1pt,
  boxsep=0pt,before skip=6pt,after skip=8pt,
  before upper={\clubpenalty=500 \widowpenalty=500 \predisplaypenalty=200 \setlist{beginpenalty=300}}}
% Cuadro complementario: por qué, cálculos, gráficos, errores y alternativas
\newtcolorbox{complemento}[1]{enhanced,breakable,colback=tceegrisclaro!45,colframe=tceedorado,
  boxrule=0pt,leftrule=3pt,arc=1pt,fontupper=\footnotesize,left=8pt,right=8pt,top=5pt,bottom=5pt,
  detach title,before upper={\clubpenalty=500 \widowpenalty=500 \setlist{beginpenalty=300}{\small\bfseries\color{tceedorado!75!black}Cuadro complementario #1}\par\smallskip},
  before skip=8pt,after skip=10pt}
% (En «examen» y «complemento», before upper relaja las penalizaciones de líneas
% viudas y huérfanas de tcee.sty, 10000: si no, la caja no encuentra dónde
% partirse y deja huecos.)
% Rótulo de una parte del texto de examen («a) …»)
% (sin \Needspace: dentro de una tcolorbox partible fuerza saltos y deja páginas casi vacías)
\newcommand{\rotulo}[1]{\par\addvspace{4pt}{\color{tceemorado}\bfseries #1}\par\nopagebreak\smallskip}
% Etiqueta dentro de un cuadro complementario
\newcommand{\que}[1]{\par\smallskip{\bfseries\color{tceemoradooscuro}#1}\enspace}
\setlength{\emergencystretch}{1.5em}

\begin{document}
\tceetitulo

% Introducción: cómo leer la solución, dónde están los retos y reparto del tiempo
\section*{Introducción}

\begin{itemize}
\item \textbf{\emph{Qué es}}: la solución \emph{de máximos} del examen de ejemplo del dictamen sobre España (10 preguntas, una por bloque de la guía, 3 horas y 30 minutos orientativas). Aplica la guía de principio a fin y remite a sus apartados («guía, I.4» es el apartado~I.4, \emph{Demanda}). Todas las cifras son del cuadernillo (cuadros 1 a 9 y 2~bis, gráficos 1 a 14, datos oficiales a 26 de octubre de 2026); lo que viene de fuera se dice y se cita.
\item \textbf{\emph{Cómo se lee}}: cada pregunta tiene 3 piezas, que se distinguen a simple vista:
  \begin{itemize}
  \item el \textbf{enunciado}, en un recuadro verde, con el \emph{peso} de la pregunta, su tiempo y la extensión del texto;
  \item el \textbf{texto de examen}, con una \textcolor{tceemorado}{\textbf{raya morada}} a la izquierda: lo que escribiría un opositor excelente, con rótulos para cada subpregunta, frase de diagnóstico al empezar y cifras citadas por su cuadro o gráfico;
  \item los \textbf{cuadros complementarios}, en gris con una \textcolor{tceedorado}{\textbf{raya dorada}}: por qué se dice cada cosa, qué apartado de la guía se aplica, los cálculos hechos con el cuadernillo, gráficos que el examen no trae, errores típicos, otras formas de redactarlo y qué añadir con más tiempo. \uline{No forman parte de la respuesta}: son la clase que hay detrás.
  \end{itemize}
\item \textbf{\emph{Un techo, no un objetivo}}: el texto de examen es más largo de lo que cabe a mano en el tiempo de cada pregunta. Está escrito para que se vea todo lo que se \emph{puede} decir; en el examen se conserva la frase de diagnóstico de cada párrafo y 1 o 2 cifras por idea (guía, D.5), y cada pregunta termina con un cuadro que dice qué recortar.
\end{itemize}

\paragraph{Dónde están los retos y cuánto pesa cada pregunta}
\begin{itemize}
\item La guía pide repartir el tiempo según los retos reales de \emph{esa} economía, no a partes iguales (guía, introducción y D.3). Leído el cuadernillo, el diagnóstico de España cabe en una frase: \uline{crece más que la zona del euro, con superávit exterior, banca rentable y deuda privada en mínimos, pero con la productividad estancada, un paro que dobla al europeo, una deuda pública en torno al 100~\% del PIB y un nuevo choque de precios en 2026}.
\item De ahí los pesos: \textbf{alto} en el perfil cíclico (ordena todo y recoge un ciclo completo con 2 recesiones), el mercado de trabajo y la productividad, la inflación, las cuentas públicas y las recomendaciones; \textbf{medio} en la inversión y el crédito, el sector exterior y la política monetaria (hay mucho que contar, pero hoy no son un problema o no se deciden en España); \textbf{bajo} en la competitividad no precio y la banca (exposición pequeña a los aranceles y banca sólida).
\item La extensión sigue a los pesos: \uline{alto~$\gg$~medio~$>$~bajo}. La tabla~\ref{tab:reparto} aplica la regla de la guía (lectura y repaso fijos; el resto, $t=W/(n+0{,}25)$ y $1{,}25\,t$ para el perfil cíclico) y desplaza minutos de las preguntas de peso bajo a las de peso alto.
\end{itemize}

\begin{table}[H]
\centering\footnotesize
\renewcommand{\arraystretch}{1.12}
\caption{Reparto orientativo del tiempo y extensión del texto de examen según el peso de cada pregunta}\label{tab:reparto}
\begin{tabular}{>{\raggedright\arraybackslash}p{4.1cm}>{\centering\arraybackslash}p{1.1cm}>{\centering\arraybackslash}p{1.25cm}>{\centering\arraybackslash}p{1.5cm}>{\raggedright\arraybackslash}p{6.4cm}}
\toprule
\rowcolor{tceeazulpalido}
\textbf{Pregunta (bloque de la guía)} & \textbf{Peso} & \textbf{Minutos} & \textbf{Palabras} & \textbf{Por qué ese peso} \\
\midrule
Lectura del cuadernillo y plan & — & \tiempo{15} & — & Fijo: fases marcadas, cálculos hechos y gráficos emparejados con su pregunta (cuadro complementario~0.1) \\
\rowcolor{tceesalvia}
1. Perfil cíclico (I) & alto & \tiempo{23} & 2.100% & Ordena todo el examen; un ciclo completo con 2 recesiones; el «porqué» del PIB por habitante \\
\rowcolor{tceesalvia}
2. Trabajo y productividad (II.1) & alto & \tiempo{21} & 1.120% & Paro estructural y productividad estancada: los 2 grandes retos de España \\
3. Consumo, inversión y crédito (II.2) & medio & \tiempo{17} & 1.408% & Buena historia (desapalancamiento, inversión más sana que en 2007), pero no es un problema; el foco, la vivienda \\
\rowcolor{tceesalvia}
4. Inflación y competitividad precio (III.1-2) & alto & \tiempo{21} & 1.579% & El problema coyuntural de 2026: 5,0~\% frente a 3,8~\%, riesgo de segunda ronda y CLU por encima de la zona \\
5. No precio y aranceles (III.3) & bajo & \tiempo{13} & 508% & Exposición directa a EE.~UU. pequeña y cuotas estables; el reto de fondo (I+D) va a la 10 \\
6. Sector exterior (IV) & medio & \tiempo{17} & 760% & La gran corrección está hecha y la PIIN mejora; quedan la energía y una PIIN aún bajo el $-35$~\% \\
7. Política monetaria (V.1) & medio & \tiempo{17} & 709% & El tono no lo decide España; sí la transmisión rápida y un tipo real más bajo que en la zona \\
8. Banca y macroprudencial (V.2) & bajo & \tiempo{13} & 432% & Banca rentable y solvente, deuda privada en mínimos; vigilar la vivienda \\
\rowcolor{tceesalvia}
9. Cuentas públicas y fiscal (V.3-4) & alto & \tiempo{21} & 1.137% & Deuda en torno al 100~\%, déficit estructural del 3~\% con brecha positiva, envejecimiento y reglas nuevas \\
\rowcolor{tceesalvia}
10. Desequilibrios y recomendaciones (VI) & alto & \tiempo{21} & 1.081% & La síntesis: donde se juntan productividad, paro, deuda y demografía \\
Repaso & — & \tiempo{11} & — & Releer, completar la frase de diagnóstico de la pregunta 1 y cerrar lo abierto \\
\midrule
\textbf{Total} & & \textbf{3 h 30 min} & \textbf{8.782}% & \\
\bottomrule
\end{tabular}
\fuente{Elaboración propia con la regla de la guía (D.3): $W=210-15-11=184$ minutos; con $n=10$, $t=184/10{,}25\approx18$ y $1{,}25\,t\approx22$. Dentro de esa media, cada pregunta de peso alto recibe 21 minutos (23 el perfil cíclico), las de peso medio 17 y las de peso bajo 13 (suman 184). Palabras: las del texto de examen de esta solución (sin enunciados ni cuadros complementarios). En verde, las preguntas de peso alto.}
\end{table}

\begin{complemento}{0.1 · Los 15 minutos de lectura: qué se marca en el cuadernillo}
\que{Qué se hace} (guía, D.4): leer las 10 preguntas enteras subrayando lo que pide cada subpregunta; mirar las unidades; dividir el ciclo con el PIB y la brecha (cuadro~1); marcar en verde y rojo fortalezas y debilidades; emparejar cada cuadro y gráfico con su pregunta, y hacer al principio los cálculos que se van a usar.

\que{Emparejar cada pieza con su pregunta.} Ningún gráfico está de adorno (guía, D.4): la tabla~\ref{tab:emparejar}, al final de este cuadro, lo hace para este examen.


\que{Los cálculos que conviene hacer al principio} (porque se usan en varias preguntas y así no hay contradicciones): el crecimiento del potencial que implica la brecha; la productividad y los CLU de 2025 ($4{,}4-0{,}0=4{,}4$); el margen ($3{,}0-4{,}4=-1{,}4$); el tipo real ($2{,}50-5{,}0=-2{,}5$); el esfuerzo fiscal ($\Delta b^{pe}=+0{,}3$ en 2025) y el saldo primario que estabiliza la deuda ($\approx-3{,}1~\%$ del PIB). Todos están desarrollados en los cuadros de cada pregunta y en el fichero \texttt{calculos.txt}.

\que{Unidades que engañan en este cuadernillo}: en el tipo de cambio efectivo real (C6, G8) un aumento es una \emph{apreciación}; la deuda externa neta va con signo positivo en C7 y negativo en G11; en la cuenta financiera, $+$ es salida neta de capital; la brecha, el potencial y los saldos ajustados son de AMECO (primavera de 2026), anterior a la última revisión del INE (PIB de 2025: 2,6~\% en C1 y 2,8~\% en AMECO).

\begin{tablacaja}
\centering\scriptsize
\renewcommand{\arraystretch}{1.08}
\caption{Qué cuadros y gráficos alimentan cada pregunta}\label{tab:emparejar}
\begin{tabular}{>{\raggedright\arraybackslash}p{3.3cm}>{\raggedright\arraybackslash}p{4.6cm}>{\raggedright\arraybackslash}p{6.3cm}}
\toprule
\rowcolor{tceeazulpalido}
\textbf{Pregunta} & \textbf{Piezas del cuadernillo} & \textbf{Para qué} \\
\midrule
1. Perfil cíclico & C1, C2, C2 bis, G1, G2 & Fases, contribuciones, oferta, 2026 trimestral, PIB por habitante \\
2. Trabajo & C3, C2 bis, G3 & Paro, Beveridge, temporalidad, productividad, CLU \\
3. Capital y crédito & C1 (FBCF), C4, G4, G5 & Ahorro, inversión por activos, deuda privada, vivienda \\
4. Inflación & C5, C3 (CLU), C6, G6, G7, G8 & Brote de 2021-2023, choque de 2026, segunda ronda, TCER \\
5. No precio & C6, G9 & Cuotas, I+D, alta tecnología, exportaciones a EE.~UU. \\
6. Sector exterior & C7, C4 (ahorro e inversión), G10, G11 & Subcuentas, financiación, PIIN y efecto valoración \\
7. Monetaria & C8, C4 (tipos), G12, G6 & Tono (tipo real), canales, reacción de 2026 \\
8. Banca & C8, C4, G4 & Semáforo de ratios, vivienda, deuda privada \\
9. Fiscal & C9, C1 (brecha), G13, G14 & Saldos, esfuerzo, $i-\gamma^N$, reglas \\
10. Recomendaciones & todo (matriz de fortalezas y debilidades) & Síntesis y riesgos de las previsiones \\
\bottomrule
\end{tabular}
\fuente{Cuadernillo del examen de ejemplo. C: cuadro; G: gráfico.}
\end{tablacaja}
\end{complemento}

% Pregunta 1. Perfil cíclico (bloque I) · peso alto
\section*{Pregunta 1. Perfil cíclico}

\begin{enunciado}
\textbf{1.} Describa brevemente el \textbf{perfil cíclico} de la economía española a lo largo del periodo 2007-2025, desde la óptica de la demanda y de la oferta, identificando sus distintas fases, y valore su situación coyuntural a la luz de los datos trimestrales de 2026 y de las previsiones para 2026 y 2027. ¿Qué indica la brecha de producción sobre la posición cíclica? ¿A qué podría deberse que, desde 2019, el PIB por habitante haya crecido bastante menos que el PIB?

\smallskip
\textbf{Peso alto} · \tiempo{23} · guía: I.2 (fases), I.3 (medidores), I.4 (demanda), I.5 (oferta), I.6 (diagnóstico) · cuadros 1, 2 y 2~bis; gráficos 1 y 2.
\end{enunciado}

\begin{examen}
\textit{La economía española ha recorrido entre 2007 y 2025 un ciclo completo con 2 recesiones y encadena desde 2021 una expansión más intensa que la de la zona del euro, apoyada en la demanda interna, el empleo y el aumento de la población. Produce ya por encima de su potencial y mantiene superávit exterior, pero crece más por acumular trabajo que por ganar productividad, y afronta en 2026 un nuevo choque de precios energéticos.}

\rotulo{a) Las fases del ciclo, por el lado de la demanda y de la oferta}
El cuadro recoge más de un ciclo: el de la crisis financiera y de la deuda soberana y el de la pandemia. Con el crecimiento del PIB real y la brecha de producción (cuadro~1) se distinguen 5 etapas:

\textbf{1.}~\textbf{Final de la expansión (2007-2008).} El PIB pasa de crecer un 3,5~\% a un 0,8~\% con la brecha todavía positiva (+3,3~\% en 2007): la economía producía por encima de su potencial, con una necesidad de financiación exterior del 9,0~\% del PIB (cuadro~7) que anticipaba un ajuste.


\textbf{2.}~\textbf{Doble recesión (2009-2013).} Una primera caída en 2009 ($-3{,}8$~\%), un estancamiento en 2010-2011 y una segunda recesión en 2012-2013 ($-2{,}9$ y $-1{,}4$~\%): una recesión en forma de W, con la brecha en un mínimo del $-9{,}6$~\% en 2013 y el paro en el 26,1~\%. La primera caída responde al fin del auge inmobiliario y del crédito; la segunda coincide con la crisis de la deuda soberana (diferencial de 435~pb en 2012, cuadro~9) y con la consolidación fiscal.


\textbf{3.}~\textbf{Recuperación y expansión (2014-2019).} El PIB crece un 2,6~\% de media; la brecha se cierra en 2017 y llega al +3,1~\% en 2019.


\textbf{4.}~\textbf{Pandemia (2020).} Una caída del 10,9~\%, mucho mayor que la de la zona del euro (en torno al 6~\%, gráfico~2), por el peso de los servicios de contacto y del turismo.


\textbf{5.}~\textbf{Rebote y nueva expansión (2021-2025).} Tras crecer un 6,7 y un 6,4~\% en 2021-2022, el PIB mantiene ritmos del 2,4-3,7~\% y la brecha es positiva desde 2022 (+2,0~\% en 2025).


Por el lado de la \textbf{demanda}, la nacional es \emph{procíclica} y la externa \emph{contracíclica}. La demanda nacional explica tanto las expansiones (4,2~pp en 2007 y en 2015; 3,4~pp en 2025) como las recesiones ($-6{,}5$~pp en 2009 y $-4{,}9$~pp en 2012), mientras la externa amortigua las caídas (+2,8~pp en 2009 y +2,0~pp en 2012), más por el desplome de las importaciones ($-18{,}3$~\% en 2009) que por el empuje de las exportaciones. Dentro de la nacional, la inversión es, como es habitual, el componente más volátil (la formación bruta de capital fijo cae un 17,1~\% en 2009 y la inversión en vivienda, un 20,2~\%), y el consumo privado sigue al PIB con menos amplitud, salvo en 2020 ($-12{,}1$~\%), cuando el cierre de la actividad lo convirtió en el canal principal del choque. El consumo público amortigua la primera recesión (+4,9 y +3,8~\% en 2008-2009) pero se recorta en plena segunda ($-3{,}9$~\% en 2012), un indicio de política fiscal procíclica que se retoma en la pregunta~9. La demanda externa vuelve a sumar en 2022 (+2,3~pp), con la vuelta del turismo, y resta en 2025 ($-0{,}8$~pp) porque las importaciones (+5,5~\%) crecen más que las exportaciones (+2,6~\%) al tirar la demanda interna: una expansión, por tanto, guiada desde dentro.

Por el lado de la \textbf{oferta} (cuadro~2), la expansión de 2014-2019 fue generalizada; la actual es sobre todo de servicios de mercado, con las actividades profesionales (+8,3 y +7,2~\% en 2024-2025) y la información y las comunicaciones a la cabeza, mientras la industria apenas avanza ($-2{,}1$~\% en 2023 y +0,6~\% en 2025), lo que sería coherente con el encarecimiento de la energía y la pérdida de competitividad en costes (pregunta~4). La agricultura es volátil y acíclica ($-16{,}9$~\% en 2022, probablemente por la sequía, y +15,0 y +11,9~\% después), y la construcción, que fue la gran perdedora de 2009-2013, repunta en 2025 (+6,0~\%).

\rotulo{b) La situación en 2026 y las previsiones}
Los datos trimestrales confirman que la expansión continúa: el PIB crece un 0,6 y un 0,7~\% en el primer y el segundo trimestre de 2026 (gráfico~1), un ritmo anualizado del 2,5-2,8~\%, apoyado en el consumo privado (0,67~pp en el segundo trimestre), mientras la demanda externa vuelve a restar ($-0{,}46$~pp). Con avances del 0,4~\% en la segunda mitad del año, el crecimiento medio de 2026 rondaría el 2,5~\%, en línea con la previsión de la Comisión (2,4~\%). Las previsiones (2,4~\% en 2026 y 1,9~\% en 2027) dibujan una \textbf{desaceleración suave}: el PIB crecería menos que su potencial (2,7 y 2,5~\%) y la brecha se reduciría al 1,7 y al 1,2~\%. Es un aterrizaje suave plausible, pero con riesgos a la baja (el petróleo, la subida de tipos del BCE y los aranceles) que se analizan en la pregunta~10.

\rotulo{c) Qué dice la brecha de producción}
La brecha es positiva desde 2022 y está en torno al 2~\% del potencial: \textbf{la economía produce por encima de su potencial}, en una fase de expansión madura. Es coherente con lo que se verá en otros bloques: una inflación por encima de la de la zona del euro, unos costes laborales unitarios que crecen el 4,4~\% y un mercado de trabajo con más vacantes que antes de la pandemia. Sin embargo, la brecha apenas se amplía desde 2024 (del 1,9 al 2,0~\%): el PIB crece ya al ritmo del potencial, que la Comisión estima en el 2,7~\%, muy por encima del $-0{,}1$-1,0~\% de 2014-2019. \uline{Un potencial tan alto para una economía avanzada no se explica por la productividad, sino por la población}: la economía crece con fuerza sin que la brecha se abra porque llegan trabajadores. Conviene leerla con prudencia: la brecha es una estimación, se revisa mucho y es especialmente incierta cuando la población activa cambia deprisa.

\rotulo{d) Por qué el PIB por habitante crece menos que el PIB}
Entre 2019 y 2025 el PIB crece un 10,0~\% y el PIB por habitante solo un 5,0~\% (gráfico~2), porque la población aumenta casi un 5~\% (cuadro~3), impulsada por la inmigración: la población extranjera pasa del 10,3 al 14,1~\% del total. La descomposición del cuadro~2~bis lo precisa: de los 1,59~puntos de crecimiento medio anual del PIB, 0,79 los aporta la población y 0,95 la mayor tasa de ocupación, mientras las horas por ocupado restan 0,46 y la productividad por hora solo suma 0,33. Es un \textbf{crecimiento extensivo}, apoyado en más personas trabajando y no en que cada hora produzca más. En la zona del euro la diferencia entre PIB y PIB por habitante es mucho menor (6,8 y 4,7~\%), y desde 2007 el PIB por habitante español ha crecido menos (8,7~\% frente al 12,6~\%): España crece más en total, pero no converge en renta por habitante. Como solo la productividad eleva la renta por habitante a largo plazo, este es el reto de fondo que recorre el resto del dictamen.

\textit{En suma, parece tratarse de una economía avanzada en una expansión madura, guiada por la demanda interna y sostenida por la población, con fortalezas (superávit exterior, desapalancamiento privado, banca saneada) y debilidades (productividad, paro estructural, deuda pública) que se desarrollan a continuación.}
\end{examen}

\begin{complemento}{1.1 · Por qué se contesta así: los 3 pasos y las fases}
\que{Qué se aplica.} Los 3 pasos del perfil cíclico (guía, I, gráfico «Los 3 pasos»): fases con el PIB y la brecha (I.2-I.3), motor por la demanda y la oferta (I.4-I.5) y diagnóstico anticipado (I.6). La frase inicial es la «frase ejecutiva» del FMI y del Banco de España (I.2): \emph{dónde está, qué la sostiene, qué la frena}; se escribe al final, en un hueco que se deja al empezar. Como referencia, el FMI resumía así su evaluación de 2026: la economía española «siguió creciendo más que la zona del euro pese al impacto adverso de la guerra en Oriente Medio» \parencite{FMI2026ArticuloIV}.

\que{Las fases, en un gráfico.} El cuadernillo no trae la brecha dibujada; el gráfico~\ref{fig:s01} la pone junto al PIB para ver las 4 fases de la guía (recuperación, expansión, desaceleración, recesión) en los datos.

\begin{figuracaja}
\caption{España: crecimiento del PIB real y brecha de producción, 2007-2027}\label{fig:s01}
\grafico[0.98\linewidth]{s01-fases-ciclo}
\fuente{Cuadro~1 del cuadernillo (Eurostat; brecha y previsiones de 2026-2027: Comisión Europea, AMECO, primavera de 2026). Qué hay que ver: la brecha toca fondo en 2013 aunque el PIB ya caía menos (\emph{pasar de caer un 2,9~\% a un 1,4~\% no es recuperación}); se cierra en 2017 y vuelve a ser positiva desde 2022; en 2026-2027 el PIB crecería menos que el potencial y la brecha bajaría: desaceleración.}
\end{figuracaja}

\begin{tablacaja}
\centering\scriptsize
\caption{Las fases en cifras (medias del periodo; brecha y paro, en el último año)}\label{tab:fases}
\begin{tabular}{>{\raggedright\arraybackslash}p{3.4cm}*{5}{>{\centering\arraybackslash}p{1.75cm}}}
\toprule
\rowcolor{tceeazulpalido}
\textbf{Fase} & \textbf{PIB (\%)} & \textbf{Demanda nacional (pp)} & \textbf{Demanda externa (pp)} & \textbf{Brecha (\% del potencial)} & \textbf{Paro (\%)} \\
\midrule
Final de la expansión, 2007-2008 & 2,2 & 1,7 & 0,5 & 1,1 & 11,3 \\
Doble recesión, 2009-2013 & $-1{,}7$ & $-3{,}5$ & 1,8 & $-9{,}6$ & 26,1 \\
Recuperación y expansión, 2014-2019 & 2,6 & 2,6 & 0,0 & 3,1 & 14,1 \\
Pandemia, 2020 & $-10{,}9$ & $-8{,}8$ & $-2{,}2$ & $-8{,}5$ & 15,5 \\
Rebote, 2021-2023 & 5,2 & 4,2 & 1,0 & 1,0 & 12,2 \\
Expansión, 2024-2025 & 3,2 & 3,4 & $-0{,}3$ & 2,0 & 10,5 \\
\bottomrule
\end{tabular}
\fuente{Cálculo propio con el cuadro~1 y el cuadro~3 (\texttt{calculos.txt}). Qué hay que ver: en las recesiones la demanda externa suma y la nacional resta; la recuperación de 2014-2019 fue de demanda nacional, con la externa neutra en promedio.}
\end{tablacaja}

\que{Errores típicos.} Contar la serie año a año («en 2010 crece un 0,1~\%, en 2011 cae un 0,6~\%…»): la guía pide la visión de conjunto, con lo atípico como excepción (I.2). Llamar «recuperación» a 2013, cuando el PIB todavía cae. Decir que la demanda interna «crece un 3,4~\%» en 2025: son 3,4 \emph{puntos} de los 2,6 que crece el PIB (I.4). Desarrollar aquí la balanza de pagos o la inflación, que tienen su pregunta: en el perfil cíclico la demanda externa se despacha en 1 o 2 frases (nota al opositor de I.4).
\end{complemento}

\begin{complemento}{1.2 · Los cálculos de la respuesta}
\que{Tasa anualizada} (guía, A.3.1): $(1+g)^4-1$; con $g=0{,}63~\%$, un 2,5~\%, y con $g=0{,}68~\%$, un 2,8~\%. Sin calculadora, $4g$ da 2,5 y 2,7: basta.

\que{Crecimiento medio de 2026 con los datos trimestrales.} Encadenando las tasas del gráfico~1 (base: cuarto trimestre de 2024 = 100), la media de 2025 es 101,5 y el segundo trimestre de 2026 está en 103,8. Si el PIB no creciera nada en la segunda mitad, 2026 crecería un 2,2~\% de media (solo por el «arrastre»); con un 0,4~\% trimestral, un 2,5~\%. La previsión de la Comisión (2,4~\%) supone, por tanto, algo menos del 0,4~\% por trimestre: una desaceleración suave, no un frenazo.

\que{Potencial que implica la brecha} (guía, I.3): $\gamma_{Y^p}\approx\gamma_Y-\Delta\mathrm{OG}$. En 2025, $2{,}6-(2{,}0-1{,}9)=2{,}5$, frente al 2,7 del cuadro; en 2024, $3{,}7-0{,}9=2{,}8$, frente al 2,5. Las diferencias se deben a que la brecha y el potencial son de AMECO (primavera de 2026), anterior a la última revisión del PIB del INE (2025: 2,8~\% en AMECO y 2,6~\% en el cuadro): \uline{si 2 fuentes no cuadran, se dice}, sin construir el análisis sobre décimas.

\que{Contribuciones} (guía, I.4): 2025, $3{,}4+(-0{,}8)=2{,}6$; 2009, $-6{,}5+2{,}8=-3{,}7\approx-3{,}8$ (los volúmenes encadenados no suman exacto).

\que{PIB por habitante} (gráfico~2): $118{,}7/107{,}9=1{,}100$ (+10,0~\%) y $108{,}7/103{,}6=1{,}050$ (+5,0~\%); población acumulada 2019-2025 (cuadro~3): $1{,}008\times1{,}005\times1{,}000\times1{,}009\times1{,}012\times1{,}011\times1{,}010\approx1{,}048$.
\end{complemento}

\begin{complemento}{1.3 · El PIB por habitante, descompuesto}
\que{La identidad} (guía, II.1, gráfico «PIB por persona en edad de trabajar»): en logaritmos,
\[
\gamma_Y=\underbrace{\gamma_N}_{\text{población}}+\underbrace{\gamma_{P_{15\text{-}64}/N}}_{\text{peso de la edad de trabajar}}+\underbrace{\gamma_{L/P_{15\text{-}64}}}_{\text{tasa de ocupación}}+\underbrace{\gamma_{H/L}}_{\text{horas por ocupado}}+\underbrace{\gamma_{Y/H}}_{\text{PIB por hora}} ,
\]
y el PIB por habitante es todo menos el primer término. El cuadro~2~bis la trae ya calculada; el gráfico~\ref{fig:s02} la dibuja por periodos.

\begin{figuracaja}
\caption{Descomposición del crecimiento del PIB por periodos (puntos al año)}\label{fig:s02}
\grafico[0.8\linewidth]{s02-descomposicion}
\fuente{Cuadro~2~bis del cuadernillo (Eurostat). Qué hay que ver: en la crisis la productividad por hora «sube» porque se destruye el empleo menos productivo (efecto composición, guía I.5); en 2013-2019 manda la ocupación; desde 2019 pesan la población y la ocupación, las horas por ocupado restan y la productividad por hora apenas suma.}
\end{figuracaja}

\que{Cómo leerla.} 3 periodos, 3 patrones: 2007-2013, ajuste por cantidades (la ocupación resta 2,9 puntos al año) con una productividad que sube de forma espuria; 2013-2019, recuperación del empleo con la población casi estancada; 2019-2025, \emph{más gente y menos horas}. La caída de las horas por ocupado ($-0{,}46$) puede deberse a más empleo a tiempo parcial, a la composición del empleo nuevo o a cambios de jornada: con el cuadernillo no se puede separar, así que se dice con condicional.

\que{Alternativa de redacción} (más corta, para el examen de 1 h 45 min): «Desde 2019 el PIB por habitante crece la mitad que el PIB (5,0 frente a 10,0~\%) porque la población aumenta casi un 5~\% por la inmigración; el crecimiento es extensivo: de 1,6 puntos anuales, 0,8 los pone la población y solo 0,3 la productividad por hora».

\que{Qué añadir con más tiempo.} La renta nacional frente al PIB (guía, I.3): las rentas primarias son casi nulas (+0,2~\% del PIB en 2025, cuadro~7), así que la RNB y el PIB apenas difieren, a diferencia de 2007-2008 ($-2{,}5$ y $-2{,}8~\%$), cuando España pagaba al exterior los intereses de su deuda. Y los indicadores adelantados (PMI, confianza), que el cuadernillo no trae: se dice que faltan.

\que{Qué recortar en el examen.} Quedarse con la frase de diagnóstico, las 5 fases en una línea cada una, 1 frase de demanda, 1 de oferta, el dato de 2026 con la previsión, 2 frases de brecha y la descomposición en 1,59 = 0,79 + 0,95 $-$ 0,46 + 0,33. Unas 500 palabras.
\end{complemento}

% Pregunta 2. Mercado de trabajo y productividad (bloque II.1) · peso alto
\clearpage
\section*{Pregunta 2. Mercado de trabajo y productividad}

\begin{enunciado}
\textbf{2.} Caracterice el \textbf{mercado laboral} español a lo largo de la serie. ¿Puede afirmarse que se encuentra tensionado? Analice la evolución de la curva de Beveridge. ¿A qué podría deberse la caída de la temporalidad desde 2022? Valore la evolución de la \textbf{productividad} y su contribución al crecimiento.

\smallskip
\textbf{Peso alto} · \tiempo{15} · guía: II.1 (cantidades, productividad, CLU y margen, tipos de paro, \autor{Okun}, \autor{Beveridge}), I.5 y V.5 · cuadros 3 y 2~bis; gráfico 3.
\end{enunciado}

\begin{examen}
\textit{El mercado de trabajo español amplifica el ciclo: destruye empleo con gran intensidad en las recesiones y lo crea con la misma rapidez en las expansiones. Hoy está en máximos de empleo y de actividad, con el paro más bajo desde 2008, pero conserva sus 2 debilidades de fondo: un paro estructural que dobla al de la zona del euro y una productividad que apenas crece.}

\rotulo{a) Caracterización a lo largo de la serie}
La tasa de paro (cuadro~3) dibuja el ciclo con nitidez: del 8,2~\% en 2007 al 26,1~\% en 2013, para bajar al 14,1~\% en 2019 y al 10,5~\% en 2025 (10,2~\% en el primer semestre de 2026, gráfico~3). Es todavía el doble que en la zona del euro (6,4~\%), con un paro juvenil del 24,9~\% (55,5~\% en 2013). La tasa de actividad está en máximos (74,9~\%) y la de empleo (67,0~\%) supera ya la de 2007: la mejora no se debe a un \emph{efecto desánimo}, sino a la creación de empleo.

Lo que más llama la atención es \textbf{cómo se ajusta}. En la crisis el ajuste fue por \emph{cantidades}: en 2009 el empleo cayó un 6,7~\% con un PIB que caía un 3,8~\%, mientras la remuneración por asalariado aún crecía un 4,0~\%; los salarios solo se moderaron después (en torno al 0~\% en 2011-2014), y los costes laborales unitarios cayeron en 2010-2014, una \emph{devaluación interna} tardía. Es la huella de un mercado dual, en el que los temporales absorbían el choque. En 2020, en cambio, el empleo cayó un 2,9~\% y las horas un 11,0~\%: el ajuste se hizo con las horas, lo que sería coherente con esquemas de mantenimiento del empleo que, en lugar de despedir, reducían la jornada.

La relación con el ciclo es intensa. En 2014-2019 el paro bajó 2,0 puntos al año con un crecimiento medio del 2,6~\%; en 2022-2025 solo 1,1 puntos con un 3,8~\%. Ahora hace falta más crecimiento para bajar el paro porque la población activa crece con la inmigración; medido frente al potencial, la respuesta del paro es parecida en las 2 fases (en torno a 1 punto de paro por punto de crecimiento por encima del potencial, el doble que el medio punto de la ley de \autor{Okun} en EE.~UU.). Y en 2023-2025 el paro bajó casi un punto al año creciendo al ritmo del potencial o algo por debajo, lo que apuntaría a una reducción del paro estructural. Ayuda que el paro de larga duración haya caído del 13,0~\% de los activos en 2013 al 3,4~\%: la crisis dejó menos \emph{histéresis} de la que cabía temer, aunque el paro, que en 2019 aún estaba en el 14~\%, sigue muy por encima de la tasa de pleno empleo de una economía avanzada.

\rotulo{b) ¿Está tensionado?}
\textbf{Algo más que antes, pero no en conjunto.} A favor de la tensión: la tasa de vacantes está en su máximo desde 2011 (0,9~\%), la relación entre vacantes y parados ha subido de 0,05 en 2019 a casi 0,09 en 2025, y la remuneración por asalariado crece un 4,4~\% con una productividad estancada, de modo que los costes laborales unitarios suben el 4,4~\% (3,3~\% en la zona del euro). En contra: una economía que produce en torno a su potencial (brecha del 0,6~\%), un paro del 10,5~\%, un paro juvenil del 25~\% y unas vacantes que siguen siendo muy pocas frente a las de la zona del euro. La tensión parece, por tanto, \emph{localizada} (sectores, ocupaciones y territorios donde falta mano de obra) más que general; el riesgo es que alimente los efectos de segunda ronda del choque energético (pregunta~4).

\rotulo{c) La curva de \autor{Beveridge}}
El gráfico~3 recorre 3 tramos. Entre 2010 y 2013 la economía baja a lo largo de la curva hacia la derecha (vacantes del 1,25 al 0,57~\%, paro del 19,9 al 26,1~\%): es el ciclo, la recesión. Entre 2014 y 2019 deshace el camino hacia arriba y a la izquierda (el paro baja al 14,1~\% y las vacantes suben al 0,70~\%): la recuperación. Tras el paréntesis de 2020, desde 2021 sigue hacia la izquierda, y desde 2022 lo hace casi en horizontal: \uline{el paro baja del 13,0 al 10,2~\% con las vacantes estables en el 0,9~\%}. Eso sería compatible con un desplazamiento de la curva hacia el origen, es decir, con un mejor emparejamiento entre parados y puestos, quizá favorecido por el cambio normativo que se comenta en el apartado~d) y por la llegada de trabajadores a los sectores con escasez. En todo caso, la curva española sigue muy alejada de la europea: muy pocas vacantes con mucho paro, la combinación de un paro estructural alto y de problemas de emparejamiento (cualificaciones, territorio, vivienda); conviene leerla por su evolución más que por su nivel, porque la medición de las vacantes también influye.

\rotulo{d) La caída de la temporalidad desde 2022}
Los asalariados temporales pasan del 25,4~\% en 2021 al 15,4~\% en 2025: 10 puntos en 4 años y \emph{en plena expansión}, cuando lo habitual en España era lo contrario (la temporalidad subió del 24,0 al 26,9~\% entre 2014 y 2018). Una caída tan brusca, persistente y a contracorriente del ciclo no parece cíclica: apunta a un \textbf{cambio normativo} que haya restringido la contratación temporal a finales de 2021. Con matices: parte del trasvase puede ser a contratos indefinidos menos estables (discontinuos o de corta duración), que la tasa de temporalidad no recoge como empleo inestable. Si se consolida, reduce la dualidad, da más estabilidad al consumo y más incentivos a formar a los trabajadores, aunque su efecto sobre la productividad está por ver.

\rotulo{e) La productividad y su contribución al crecimiento}
La productividad es \textbf{el punto débil}. Por hora trabajada crece solo un 0,33~\% anual entre 2019 y 2025 (1,9~\% acumulado) y por ocupado cae ligeramente ($-0{,}13~\%$ anual), porque se trabajan menos horas por persona (cuadros~2~bis y~3). De los 1,59~puntos de crecimiento medio del PIB, la productividad por hora aporta 0,33, una quinta parte: el resto lo pone el trabajo. En 2007-2013 la productividad creció más (1,6~\% anual), pero de forma espuria, porque se destruyó sobre todo el empleo menos productivo (construcción, temporales): en España la productividad aparente suele ser \emph{contracíclica}.

Como $\gamma_{\mathrm{CLU}}=\gamma_W-\gamma_{\mathrm{PAT}}$, con una productividad estancada todo el aumento salarial pasa a los costes: en 2025, $4{,}4-0{,}0=4{,}4~\%$. Las causas de fondo serían el peso de los sectores de baja productividad, el pequeño tamaño medio de las empresas, el escaso esfuerzo en I+D (1,50~\% del PIB, frente al 2,26~\% de la UE, cuadro~6) y una inversión en maquinaria que apenas crece desde 2019. Con una población que envejece (la tasa de dependencia de los mayores pasa del 24,0 al 31,2~\%), \uline{el crecimiento futuro dependerá cada vez más de la productividad y no de sumar trabajadores}: es el reto que justifica las reformas de la pregunta~7.
\end{examen}

\begin{complemento}{2.1 · Por qué se contesta así: el tipo de paro decide la receta}
\que{Qué se aplica.} Guía, II.1: los rótulos con cada subpregunta al principio (nota al opositor de II.1: «para que el tribunal las encuentre al oírlas»); la tasa de actividad para descartar el efecto desánimo; la flexibilidad del mercado (¿se ajusta por precios o por cantidades?); el tipo de paro (gráfico~\ref{fig:tipos-paro}); \autor{Okun} con $\Delta U$; \autor{Beveridge}; y la cadena productividad-CLU-margen.

\que{Un dato de fuera que encaja.} La Comisión estima una NAWRU de España en torno al 12~\% \parencite{CE2026AMECO} (guía, II.1). Con un paro del 10,5~\%, la economía estaría por debajo de su paro de equilibrio, lo que cuadraría con unos CLU que se aceleran; el FMI, en cambio, sitúa el paro estructural en torno al 10~\%, es decir, que el paro actual ya estaría en su nivel estructural \fmi{6}, coherente con la brecha casi nula del cuadernillo. Es la respuesta de máximos a «¿está tensionado?»: \emph{para su propio paro estructural, sí; para el de la zona del euro, no}.

\begin{figuracaja}
\caption{Tipos de paro: cómo se detectan y qué receta tienen (infografía de la guía)}\label{fig:tipos-paro}
\grafico[0.52\linewidth]{esquema-tipos-paro}
\fuente{Guía del dictamen, II.1. Qué hay que ver: cada tipo deja una huella distinta en el cuadro; la tabla siguiente la busca en el cuadernillo de España.}
\end{figuracaja}

\begin{tablacaja}
\centering\scriptsize
\caption{¿Qué paro tiene España? Los indicios del cuadernillo}\label{tab:tipo-paro}
\begin{tabular}{>{\raggedright\arraybackslash}p{2.6cm}>{\raggedright\arraybackslash}p{6.2cm}>{\raggedright\arraybackslash}p{5.2cm}}
\toprule
\rowcolor{tceeazulpalido}
\textbf{Tipo de paro} & \textbf{Indicio en el cuadernillo} & \textbf{Lectura} \\
\midrule
Keynesiano (demanda) & Consumo privado +3,1~\% y brecha +0,6~\% en 2025 & \textcolor{tceeverde}{No}: la demanda tira \\
Clásico (salarios) & CLU +4,4~\% (zona del euro 3,3~\%); CLU reales +1,4 en 2025; sin embargo, de 2007 a 2025 los CLU crecen menos que en la zona & \textcolor{tceedorado}{Riesgo reciente}, no un problema acumulado \\
Friccional o de emparejamiento & Vacantes muy bajas con paro alto; paro juvenil del 24,9~\%; Beveridge lejos de la europea & \textcolor{tceerojo}{Sí}: el problema principal \\
De larga duración & 3,4~\% de los activos (13,0~\% en 2013) & \textcolor{tceeverde}{Mejora mucho}; histéresis menor \\
\bottomrule
\end{tabular}
\fuente{Cuadros 1, 3 y 6 y gráfico 3 del cuadernillo. Recetas: guía, V.5 (políticas activas, formación, movilidad y vivienda para el emparejamiento).}
\end{tablacaja}

\end{complemento}

\begin{complemento}{2.2 · La ley de \autor{Okun} con los datos}
\que{La fórmula de la guía} (II.1): $\Delta U_t\approx-0{,}5\,(\gamma_{Y,t}-\gamma_{Y^p,t})$, con $U$ la tasa de paro. Se aplica con el potencial del cuadro~1, no con el crecimiento a secas.


\que{Cálculo}: 2014-2019, $\Delta U$ medio $=-2{,}0$; $\gamma_Y-\gamma_{Y^p}=2{,}62-0{,}96=1{,}7$; coeficiente $-2{,}0/1{,}7\approx-1{,}2$. 2022-2025, $\Delta U=-1{,}1$; $\gamma_Y-\gamma_{Y^p}=3{,}75-2{,}59=1{,}2$; coeficiente $\approx-0{,}95$. \uline{El coeficiente casi no cambia: lo que cambia es el potencial}, que ahora incluye la llegada de activos.

\que{Por qué baja el paro con crecimiento igual al potencial.} Si el potencial del FMI incorpora una NAWRU que baja, el paro observado baja con ella aunque la brecha no cambie. Es la lectura optimista (reforma laboral, mejor emparejamiento); la prudente, que el potencial esté sobrestimado. Las 2 se dicen con condicional.

\que{Error típico.} Decir que «la ley de \autor{Okun} no se cumple» porque en 2022-2025 el paro baja menos con más crecimiento: la guía advierte que el coeficiente varía y que hay que comparar con el potencial (II.1).

\begin{figuracaja}
\caption{España: variación del paro y crecimiento por encima del potencial, 2008-2025}\label{fig:s03}
\grafico[0.74\linewidth]{s03-okun}
\fuente{Cuadros~1 y~3 del cuadernillo (FMI, WEO de octubre de 2026: PIB, potencial y paro). Qué hay que ver: los puntos de 2014-2019 caen por debajo de la recta de \autor{Okun} de EE.~UU.: en España el paro reacciona más al crecimiento; en 2023-2025 (crecimiento igual o algo inferior al potencial) el paro baja casi un punto al año; 2020-2021 se salen de la relación (ERTE).}
\end{figuracaja}
\end{complemento}

\begin{complemento}{2.3 · \autor{Beveridge}, temporalidad y productividad: cálculos y contexto}
\que{Tensión: vacantes por parado} (guía, II.1: «la ratio vacantes/parados resume la tensión»): $0{,}70/14{,}1=0{,}050$ en 2019; $0{,}90/10{,}5=0{,}086$ en 2025; $0{,}90/10{,}25=0{,}088$ en el primer semestre de 2026. En la zona del euro, $2{,}25/6{,}3=0{,}36$ en 2025, 4 veces más \parencite{Eurostat2026jvs} (dato de la guía, no del cuadernillo).

\que{La serie de vacantes tiene un salto} entre 2009 y 2010 (el cuadernillo la da desde 2010): no se compara 2009 con 2010.

\que{Temporalidad y ciclo.} Antes de 2022 era procíclica: bajó en la crisis (del 31,6 al 23,2~\% entre 2007 y 2013, porque se destruía sobre todo empleo temporal) y subió en la recuperación. Su caída en 2022-2025, con el PIB creciendo, rompe el patrón: \emph{un cambio de patrón que coincide con un cambio normativo} es la prueba indirecta que el tribunal espera.

\que{Productividad, en 3 periodos} (cuadro~3, acumulada): por hora, +10,2~\% en 2007-2013, +2,7~\% en 2013-2019 y +1,9~\% en 2019-2025; por ocupado, +9,8, +2,2 y $-0{,}8~\%$. \uline{La que más crece es la de la crisis}: efecto composición (guía, I.5, «cuidado»).

\que{Salario real y productividad, 2019-2025}: la remuneración por asalariado sube un 25,8~\% y el IPCA un 21,5~\%: salario real de consumo +3,5~\%, por encima de la productividad por ocupado ($-0{,}8~\%$) y por hora (+1,9~\%). Es otra forma de ver por qué los CLU reales suben en 2025 y los márgenes se estrechan (pregunta~4).

\que{Alternativas de redacción.} Para la tensión: «el mercado está más tensionado que en 2019, pero lejos de la tensión europea: la escasez de mano de obra convive con un paro alto, la marca de un problema de emparejamiento». Para la temporalidad, si no se quiere nombrar la reforma: «un cambio institucional que limita la contratación temporal».

\que{Qué añadir con más tiempo.} La demografía (guía, B.1): el envejecimiento resta crecimiento potencial y la inmigración lo suma, con efectos sobre el paro a corto plazo (la inmigración eleva primero el paro y después el potencial). Y el salario mínimo, que el cuadernillo no trae: si se menciona, sin cifras y con condicional.

\que{Qué recortar en el examen.} La frase de diagnóstico; el paro en 3 cifras (2007, 2013, 2025) frente a la zona; el ajuste por cantidades en 1 frase; tensión en 2 frases; \autor{Beveridge} en 3 tramos; temporalidad en 2 frases; productividad con la cifra de 0,33 y el CLU de 2025.
\end{complemento}

\begin{complemento}{2.4 · Contexto y fuentes: lo que el cuadernillo no trae}
\que{La reforma laboral.} El cambio normativo que el texto de examen deduce de los datos es el Real Decreto-ley 32/2021, de 28 de diciembre, de medidas urgentes para la reforma laboral: limitó los contratos temporales (causas tasadas, duraciones máximas) y reforzó el contrato fijo discontinuo. Por eso la tasa de temporalidad cae en 2022, el primer año completo de aplicación. Nombrarla en el examen es razonable porque el país está identificado, pero la prueba que puntúa es la de los datos: \emph{una caída brusca, persistente y a contracorriente del ciclo}. El debate abierto (que se dice con condicional) es cuánto empleo pasó a fijos discontinuos, que la EPA cuenta como indefinidos.

\que{Los esquemas de mantenimiento del empleo.} Los expedientes de regulación temporal de empleo (ERTE) de la pandemia (Real Decreto-ley 8/2020, de 17 de marzo) permitieron suspender contratos o reducir jornada con prestación pública: explican que en 2020 el empleo cayera un 2,9~\% y las horas un 11,0~\% (cuadro~3), y que la guía (II.1) advierta que en la pandemia el paro apenas reflejó la caída del PIB.

\que{La NAWRU y la curva europea.} La NAWRU de la Comisión para España (en torno al 12~\%, \cite{CE2026AMECO}) y la curva de \autor{Beveridge} de la zona del euro (vacantes del 2,25~\% con un paro del 6,3~\% en 2025, \cite{Eurostat2026jvs}) están en los cuadros 2.1 y 2.3; no en el texto de examen, porque no vienen en el cuadernillo.

\que{Las vacantes, con cautela.} La serie de vacantes de Eurostat para España salta entre 2009 y 2010 sin marca de ruptura (por eso el cuadernillo la da desde 2010): se leen sus cambios, no su nivel frente a otros países.
\end{complemento}

\begin{complemento}{2.5 · El diagnóstico del FMI (artículo~IV de 2026) frente a los datos}
\que{Paro estructural.} El paro «parece estabilizarse en un nivel estructural de en torno al 10~\%» \fmi{6}; el paro estructural habría bajado en la última década, como sugieren el desplazamiento de la curva de \autor{Beveridge} hacia la izquierda y la ausencia de aceleración de los salarios reales pese a la caída continua del paro \fmi{27}. \textbf{Coincide} con la lectura de la solución (curva que se desplaza hacia el origen, paro estructural que baja) y aporta una cifra que el cuadernillo no da. Con un paro del 10,5~\%, la economía estaría casi en su paro estructural, coherente con la brecha casi nula.

\que{Salarios.} Los salarios pactados en convenio crecían en torno al 3,5~\% a finales de 2025 y los efectivos por asalariado, un 4,9~\%: la diferencia (\emph{deriva salarial}) refleja un mercado tensionado y el intento de recuperar la inflación pasada \fmi{7}. Es la base del riesgo de segunda ronda de la pregunta~4.

\que{La reforma de 2021.} Según el FMI, redujo de forma duradera la temporalidad y mejoró la estabilidad del empleo, aunque «modestamente», porque bajó la duración de los contratos indefinidos y subió el peso de los fijos discontinuos \fmi{27}: \textbf{coincide} con el matiz del texto de examen.

\que{Recomendaciones del FMI} \fmi{28-29}: cotizaciones más altas para las empresas que despiden y contratan con más frecuencia (\emph{experience rating}); financiar los servicios públicos de empleo según su tasa de colocación; ampliar más allá de 6 meses la compatibilidad del subsidio con el empleo; no subir el salario mínimo por encima del objetivo del 60~\% del salario medio neto ni indexarlo al IPC; y un crédito fiscal por trabajo para los hogares de renta baja. El FMI subraya además que la tasa de empleo sigue siendo de las más bajas de la zona del euro \fmicp{1-2}.

\que{Dónde difiere.} El FMI duda de que el paro pueda bajar mucho más sin nuevas reformas que no están sobre la mesa \fmi{27}; la solución es algo más optimista al leer la caída del paro de 2023-2025 con un crecimiento igual al potencial como reducción del paro estructural. Las 2 lecturas son compatibles: el paro estructural bajó, pero hasta el 10~\%.
\end{complemento}

% Pregunta 3. Consumo, ahorro, inversión y crédito (bloque II.2) · peso medio
\section*{Pregunta 3. Consumo, ahorro, inversión y crédito}

\begin{enunciado}
\textbf{3.} Analice el comportamiento del \textbf{consumo, el ahorro y la inversión} del sector privado, así como sus principales determinantes. ¿Qué factores podrían explicar el dinamismo de la inversión en los últimos años y su distinta composición respecto de la fase expansiva anterior a 2008? Valore la evolución del \textbf{endeudamiento y del crédito} al sector privado.

\smallskip
\textbf{Peso medio} · \tiempo{17} · guía: I.4 (consumo e inversión: determinantes), II.2 (crédito), IV.1 (ahorro e inversión), V.2 (endeudamiento, vivienda) · cuadros 1 y 4; gráficos 4 y 5.
\end{enunciado}

\begin{examen}
\textit{El sector privado español ha pasado de gastar a crédito, como en 2007, a ahorrar más de lo que invierte: los hogares ahorran el doble que antes de la crisis, la inversión crece con otra composición (intangibles e infraestructuras, no vivienda) y la deuda privada se ha reducido a la mitad. El único foco de riesgo es el precio de la vivienda.}

\rotulo{a) Consumo, ahorro e inversión y sus determinantes}
La tasa de ahorro de los hogares (cuadro~4, gráfico~5) resume los episodios del ciclo: sube del 6,0 al 11,7~\% entre 2007 y 2009 por \emph{precaución}; baja al 7,0~\% en 2012, cuando la renta disponible real por habitante cae un 6,6~\% y los hogares desahorran para sostener el consumo; se dispara al 17,6~\% en 2020 por el ahorro \emph{forzoso} y el precautorio; y se gasta en 2021-2022 (9,0~\%), cuando la inflación erosiona la renta real ($-2{,}0~\%$ en 2022). Desde 2023 la renta real se recupera (+3,2 y +3,3~\% por habitante en 2023-2024) más deprisa que el consumo, y el ahorro sube al 12,8~\%. En 2025 ocurre lo contrario: el consumo crece un 3,1~\% con una renta por habitante que solo avanza un 0,9~\%, y la tasa de ahorro baja al 11,9~\%. Entre los determinantes habituales, el consumo se apoya hoy en el \textbf{empleo} (+2,6~\%) y en la \textbf{población}, en la riqueza inmobiliaria y en unos tipos que bajaron en 2025; frena la renta real por persona. Que el ahorro doble el de 2007 sería coherente con el desapalancamiento, la incertidumbre y el envejecimiento.

La tasa de inversión de las empresas (24,4~\% del valor añadido) está lejos del 30,2~\% de 2007, y la de los hogares en vivienda (7,4~\% de su renta), muy lejos del 15,8~\%. Los determinantes son mixtos: la \textbf{brecha positiva} y la demanda esperada animan a ampliar capacidad (acelerador); el \textbf{coste de la financiación} subió con fuerza en 2023-2024 (préstamos nuevos a empresas al 4,6-4,7~\%) y baja en 2025 (3,4~\%, apenas 0,4 puntos por encima del deflactor del PIB); la \textbf{rentabilidad} empresarial se estrecha (los beneficios pasan del 43,8~\% del valor añadido en 2016 al 38,9~\%); y la \textbf{incertidumbre} (energía, aranceles) juega en contra.

\rotulo{b) El dinamismo de la inversión y su composición}
La formación bruta de capital fijo crece un 13,4~\% entre 2019 y 2025, pero con una composición muy distinta de la de antes de 2008: los \textbf{productos de la propiedad intelectual} (programas informáticos, I+D, bases de datos) aumentan un 47,4~\% y las \textbf{otras construcciones} (infraestructuras) un 14,9~\%, mientras la vivienda solo crece un 6,0~\% y la maquinaria un 2,7~\% (cuadro~1). Antes de 2008 tiraba la vivienda, que se hundió un 51,6~\% entre 2007 y 2013. Detrás del cambio habría 3 factores: la \textbf{digitalización}, que desplaza la inversión hacia los intangibles; los \textbf{fondos europeos}, que llegan como transferencias de capital (la cuenta de capital ronda el 1~\% del PIB desde 2021, cuadro~7) y financian infraestructuras, digitalización y transición energética; y una \textbf{oferta de vivienda rígida} (suelo, plazos, menos crédito promotor), que explica que la inversión residencial apenas crezca pese a la demanda. Es una inversión más sana que la de 2007, financiada con ahorro nacional (la tasa de ahorro nacional sube al 25,0~\% del PIB y supera a la inversión, del 21,7~\%) y que eleva el potencial, porque \uline{la inversión es demanda hoy y oferta mañana}.

\rotulo{c) Endeudamiento y crédito}
El \textbf{desapalancamiento} ha sido profundo y ordenado. El flujo de crédito al sector privado pasa del 26,5~\% del PIB en 2007 al $-11{,}2$~\% en 2012 y al 2,6~\% en 2025 (cuadro~4); la deuda de los hogares baja del 85,0~\% del PIB en 2009 al 42,9~\%, y la de las empresas, del 119,0 al 62,0~\% (gráfico~4), ambas por debajo de los umbrales del procedimiento de desequilibrios (55 y 85~\%). Desde 2019 el ajuste ha sido sobre todo por el denominador: la deuda de los hogares apenas varía en euros y su ratio baja porque el PIB nominal crece un 35~\%. Es un \emph{desapalancamiento sano}, hecho con más renta y no hundiendo el gasto. El crédito, que cayó en 2023 ($-2{,}0~\%$ a hogares y $-3{,}7~\%$ a empresas) con la subida de tipos, se reactiva en 2025 (+4,5 y +4,2~\%): la política monetaria se transmite y no hay racionamiento.

La excepción es la \textbf{vivienda}: su precio sube un 12,7~\% en 2025 (9,7~\% en términos reales), por encima del umbral del 9~\% del procedimiento de desequilibrios, aunque en términos reales sigue un 13~\% por debajo de 2007 (gráfico~4). A diferencia de entonces, no la impulsa el crédito, sino una demanda que crece con la población frente a una oferta que no responde: es más un problema de \emph{accesibilidad} que de estabilidad financiera, aunque conviene vigilarlo (pregunta~8).
\end{examen}

\begin{complemento}{3.1 · Por qué se contesta así: los determinantes, con su signo y su dato}
\que{Qué se aplica.} Guía, I.4: las funciones de consumo e inversión con sus signos, eligiendo los 2 o 3 determinantes que el cuadro permite ver y diciendo cuáles faltan (nota al opositor de I.4); II.2 (crédito como factor); IV.1 ($S-I$); V.2 (desapalancamiento sano, vivienda).

\begin{tablacaja}
\centering\scriptsize
\caption{Los determinantes del consumo y la inversión en el cuadernillo (2023-2025)}\label{tab:determinantes}
\begin{tabular}{>{\raggedright\arraybackslash}p{3.1cm}>{\centering\arraybackslash}p{0.9cm}>{\raggedright\arraybackslash}p{6.0cm}>{\raggedright\arraybackslash}p{3.6cm}}
\toprule
\rowcolor{tceeazulpalido}
\textbf{Determinante} & \textbf{Signo} & \textbf{Dato} & \textbf{Efecto} \\
\midrule
\rowcolor{tceesalvia}\multicolumn{4}{l}{\emph{Consumo}: $C=f(Y^d_{\text{real}}{}^{+},\ \text{riqueza}^{+},\ r^{-},\ \pi^{-},\ \text{crédito}^{+},\ \text{deuda}^{-},\ N^{+},\ \text{incertidumbre}^{-})$} \\
Renta disponible real & $+$ & +3,2, +3,3 y +0,9~\% por habitante (C4) & Sostiene 2023-2024; frena en 2025 \\
Empleo y población & $+$ & Empleo +2,6~\%, población +1,0~\% en 2025 (C3) & Tira del consumo total \\
Riqueza (vivienda) & $+$ & Precio +12,7~\% en 2025 (C4) & Empuja (efecto riqueza) \\
Tipo real & $-$ & Hipotecas al 2,73~\% con IPCA del 2,7~\%: tipo real $\approx0$ (C4, C5) & Neutro o favorable \\
Deuda de los hogares & $-$ & 42,9~\% del PIB, mínimo de la serie (C4) & Ya no frena \\
\rowcolor{tceesalvia}\multicolumn{4}{l}{\emph{Inversión}: $I=f(r^{lp\,-},\ \gamma^e_Y{}^{+},\ \mathrm{OG}^{+},\ \mathrm{EBE}^{+},\ \text{crédito}^{+},\ \text{incertidumbre}^{-})$} \\
Brecha y demanda & $+$ & Brecha +2,0~\%; PIB +2,6~\% (C1) & Empuja (acelerador) \\
Coste de financiación & $-$ & Empresas: 4,64~\% (2023) → 3,40~\% (2025); real $\approx3{,}40-3{,}0=0{,}4$ & Frena en 2023-2024, alivia en 2025 \\
Rentabilidad & $+$ & Beneficios 40,3~\% del VAB (2022) → 38,9~\% (C4); margen $-1{,}4$ en 2025 & Frena \\
Fondos europeos & $+$ & Cuenta de capital 0,9-1,1~\% del PIB desde 2021 (C7) & Empuja (infraestructuras, intangibles) \\
Incertidumbre & $-$ & No hay indicador; energía y aranceles (G7, G9) & Se dice con condicional \\
\bottomrule
\end{tabular}
\fuente{Cuadernillo del examen (C: cuadro; G: gráfico). Signos: guía, I.4. El tipo real a largo se aproxima con el tipo de los préstamos nuevos menos el deflactor del PIB (la guía, I.4, nota, da el orden de preferencia).}
\end{tablacaja}

\que{Error típico.} Explicar la caída del ahorro de 2025 con «los hogares consumen más porque tienen más renta»: es al revés, el consumo crece \emph{más} que la renta por habitante; lo explican el empleo, la población y la riqueza. Y decir que el consumo crece «un 3,1~\% por habitante»: el 3,1 es total; por habitante, en torno al 2~\% ($3{,}1-1{,}0$).
\end{complemento}

\begin{complemento}{3.2 · Los cálculos: inversión por activos, ahorro e inversión y deuda}
\que{Inversión acumulada por activos} (encadenando las tasas del cuadro~1):

\begin{tablacaja}
\centering\scriptsize
\caption{Formación bruta de capital fijo por activos: variación acumulada en volumen (\%)}\label{tab:fbcf}
\begin{tabular}{>{\raggedright\arraybackslash}p{4.6cm}*{4}{>{\centering\arraybackslash}p{1.8cm}}}
\toprule
\rowcolor{tceeazulpalido}
\textbf{Activo} & \textbf{2007-2013} & \textbf{2013-2019} & \textbf{2019-2025} & \textbf{2007-2025} \\
\midrule
Formación bruta de capital fijo & $-37{,}4$ & 33,8 & 13,4 & $-5{,}0$ \\
\hspace*{0.8em}Viviendas & $-51{,}6$ & 61,9 & 6,0 & $-16{,}9$ \\
\hspace*{0.8em}Otras construcciones & $-44{,}6$ & 14,8 & 14,9 & $-27{,}0$ \\
\hspace*{0.8em}Maquinaria, bienes de equipo y armamento & $-28{,}7$ & 36,3 & 2,7 & $-0{,}2$ \\
\hspace*{0.8em}Productos de la propiedad intelectual & 19,7 & 21,7 & 47,4 & 114,7 \\
PIB & $-7{,}6$ & 16,9 & 10,2 & 19,0 \\
\bottomrule
\end{tabular}
\fuente{Cálculo propio con el cuadro~1 ($\prod_t(1+\gamma_t)-1$; \texttt{calculos.txt}). Qué hay que ver: la inversión total aún no ha recuperado el nivel de 2007; solo los intangibles lo superan (más que lo duplican).}
\end{tablacaja}

\que{Ahorro, inversión y saldo exterior} (guía, IV.1 y A.3.4): $S/\mathrm{PIB}=I/\mathrm{PIB}+\mathrm{CC}/\mathrm{PIB}$. En 2025, $25{,}0-21{,}7=3{,}3\approx3{,}2$ (cuenta corriente, cuadro~7); en 2007, $21{,}1-30{,}5=-9{,}4$, exactamente el déficit corriente. El paso de un $-9{,}4$ a un +3,2 vino primero de la inversión (cae 13,1 puntos entre 2007 y 2013) y después del ahorro (sube 5,5 puntos entre 2013 y 2025): enlaza con la pregunta~6.

\que{Desapalancamiento: numerador o denominador.} Deuda en euros $=$ ratio $\times$ PIB nominal (cuadro~1). Hogares: $0{,}565\times1.254\approx708$ miles de millones en 2019 y $0{,}429\times1.690\approx725$ en 2025 (+2~\%); si la deuda no hubiera variado en euros, la ratio habría bajado sola al 41,9~\%. Entre 2009 y 2019, en cambio, la deuda en euros cayó un 22~\%: \uline{primero un desapalancamiento activo (amortizar), después uno pasivo (crecer)}. Empresas: +10,6~\% en euros entre 2019 y 2025 (los préstamos con aval público de 2020: préstamos a empresas +9,1~\% ese año).

\que{Vivienda} (cuadro~4): precio nominal acumulado 2013-2025, +88~\%; real, +48~\%; real 2019-2025, +19~\% con una inversión en vivienda que solo crece un 6,0~\%: \emph{más precio que cantidad}, la señal que la guía (V.2) asocia a tensiones.

\que{Qué añadir con más tiempo.} La relación entre precio de la vivienda y renta (no viene en el cuadernillo) y el esfuerzo de acceso (umbral orientativo de un tercio de la renta, guía V.2). \que{Qué recortar en el examen.} Las tasas de inversión de hogares y empresas (basta la composición por activos) y la descomposición del desapalancamiento (basta «la ratio baja sobre todo porque crece el PIB nominal»).
\end{complemento}

% Pregunta 4. Inflación y competitividad precio (bloque III.1-III.2) · peso alto
\section*{Pregunta 4. Inflación y competitividad precio}

\begin{enunciado}
\textbf{4.} Analice la evolución de la \textbf{inflación} desde 2021 y sus principales determinantes. ¿Qué factores podrían explicar que en 2026 la inflación española se sitúe por encima de la de la zona del euro? ¿Aprecia riesgo de efectos de segunda vuelta? Valore la evolución de la \textbf{competitividad precio y coste} a lo largo de la serie.

\smallskip
\textbf{Peso alto} · \tiempo{21} · guía: III.1 (índices, subyacente, efecto base, origen, segunda ronda), III.2 (tipo de cambio real), II.1 (CLU y margen), B.1 (un brote de inflación) · cuadros 3, 5 y 6; gráficos 6, 7 y 8.
\end{enunciado}

\begin{examen}
\textit{España ha sufrido 2 choques de precios de origen energético: el de 2021-2023, del que salió antes que la zona del euro, y el de 2026, en el que su inflación (5,0~\% en septiembre) supera a la europea (3,8~\%) sobre todo porque su inflación subyacente es más alta. Hay riesgo de efectos de segunda ronda, y la competitividad en costes ganada en 2008-2018 se ha erosionado en parte.}

\rotulo{a) La inflación desde 2021 y sus determinantes}
El brote de 2021-2023 fue, en su origen, un \textbf{choque de oferta importado}: la energía sube un 21,1~\% en 2021 y un 27,6~\% en 2022 y el deflactor de las importaciones, un 22,4~\% en 2022 (cuadro~5). Que el IPCA (8,3~\%) creciera mucho más que el deflactor del PIB (4,7~\%) confirma que la inflación venía de fuera, porque el IPC incluye lo importado y el deflactor no. El choque se trasladó después a los alimentos (+10,7 y +11,1~\% en 2022-2023) y a la subyacente (3,8 y 4,1~\%), los \emph{efectos indirectos}. Hubo también un componente de demanda: la reapertura y el ahorro embalsado en 2021-2022, con la brecha ya positiva desde 2022. Lo que \emph{no} hubo fue una espiral de precios y salarios: en 2022 la remuneración por asalariado creció un 4,7~\%, casi 4 puntos menos que los precios, de modo que los trabajadores absorbieron buena parte del choque (el salario real cayó un 3,6~\%), mientras los márgenes empresariales se ampliaban (el deflactor del PIB creció 2,6 puntos más que los CLU). Como importador neto de energía, España se empobreció frente al exterior: la relación real de intercambio empeoró en torno a un 6~\% en 2022.

España llegó antes al pico (10,7~\% en julio de 2022, frente al 10,6~\% de la zona del euro en octubre, gráfico~6) y bajó también antes: un 3,4~\% de media en 2023, frente al 5,4~\%, gracias a la caída de la energía ($-16{,}1$~\%). Lo explicarían un traslado más rápido de los precios mayoristas de la electricidad a los consumidores, en los 2 sentidos, y medidas como el tope al precio del gas. El \emph{efecto base} se ve con claridad en marzo de 2023: la inflación española cae del 6,0 al 3,1~\% en un mes porque sale de la comparación la subida de marzo de 2022. La subyacente, en cambio, tocó techo más tarde (5,1~\% a principios de 2023) y bajó despacio, por los servicios: la «última milla». En 2024-2025 la inflación española (2,9 y 2,7~\%) quedó algo por encima de la europea (2,4 y 2,1~\%), con una subyacente media del 2,8 y el 2,6~\%: la economía española llegó al choque de 2026 sin haber completado del todo la desinflación.

\rotulo{b) Por qué en 2026 la inflación española supera a la europea}
El nuevo choque es otra vez energético: el Brent pasa de 62,5 dólares en diciembre de 2025 a 117,3 en abril y 114,2 en septiembre de 2026 (gráfico~7), una subida de más del 80~\%, algo mayor en euros. La inflación española sube del 2,4~\% en enero al 5,0~\% en septiembre, y la europea, del 1,7 al 3,8~\%. Lo revelador es \textbf{dónde está la diferencia}: de los 1,2 puntos de diferencia en septiembre, 1,0 están en la subyacente (3,5~\% frente a 2,5~\%), mientras la distancia entre la inflación general y la subyacente, que aproxima el peso de la energía y los alimentos, es parecida (1,5 y 1,3 puntos). La diferencia, por tanto, \uline{no viene tanto del petróleo como de factores internos}:
\begin{itemize}
\item una \textbf{demanda más fuerte}: brecha positiva del 2~\% y un consumo que sigue tirando en 2026 (pregunta~1);
\item \textbf{costes laborales} que crecen más: CLU del 4,4~\% en 2025, frente al 3,3~\% en la zona (cuadro~3);
\item el peso de los \textbf{servicios} (turismo, hostelería) y de la \textbf{vivienda}, cuyo precio sube a 2 dígitos (pregunta~3);
\item y, en menor medida, un traslado más rápido de la energía a la factura eléctrica.
\end{itemize}
La previsión del 3,0~\% para 2026 ha quedado desfasada: la media de enero a septiembre ya es del 3,6~\%.

\rotulo{c) ¿Riesgo de efectos de segunda vuelta?}
\textbf{Sí, y mayor que en la zona del euro.} Los indicios: la subyacente española sube del 3,0~\% en enero al 3,6-3,7~\% en mayo-julio y se mantiene en el 3,5~\%, mientras la europea sigue en el 2,2-2,6~\%; los salarios ya crecían el 4,4~\% en 2025 con la productividad estancada; y la brecha y la tensión del mercado de trabajo dan poder de negociación. Además, las empresas llegan al choque con los márgenes estrechados: el deflactor del PIB creció menos que los CLU en 2025 ($\gamma_{1+\mu}\approx3{,}0-4{,}4=-1{,}4$~puntos), así que tenderán a trasladar el encarecimiento de la energía para recuperarlos. Si los salarios se indexan a la inflación pasada, se arma la espiral. Lo atenúan unas expectativas que, a juzgar por la subyacente europea, siguen ancladas, y la subida de tipos del BCE (pregunta~7). En la terminología de la guía, la inflación europea tiene rasgos de \emph{transitoria} (concentrada en la energía, con la subyacente contenida), y la española, de \emph{persistente} (subyacente al alza, salarios acelerados y brecha positiva). Con una política monetaria común que no puede responder a España, la respuesta adecuada sería un \textbf{pacto de rentas} que reparta la pérdida de renta real del choque energético entre salarios y beneficios, una política fiscal que no añada demanda (pregunta~9) y medidas que aumenten la competencia en los mercados de bienes y servicios, para que el choque no se traslade con margen.

\rotulo{d) La competitividad precio y coste a lo largo de la serie}
La serie tiene 3 tiempos (cuadros~3 y~6, gráfico~8). Hasta 2008, \textbf{pérdida}: el tipo de cambio efectivo real con CLU frente a la zona del euro (un aumento es una apreciación) sube a 111,9. Entre 2008 y 2018, \textbf{devaluación interna}: sin tipo de cambio propio, la competitividad se recupera por los costes; los CLU crecen un 1,2~\% acumulado en 2007-2013, frente al 12,1~\% en la zona, y el tipo real cae casi un 16~\% hasta 94,4 en 2018. Desde 2019, \textbf{erosión parcial}: los CLU suben un 26,8~\% en 2019-2025, frente al 23,0~\% en la zona, y el tipo real se aprecia casi un 5~\% desde su mínimo de 2018. En el conjunto de 2007 a 2025 España conserva una ganancia de competitividad en costes (el tipo real con CLU cae un 9~\%), mientras que con precios de consumo frente a 42 socios apenas varía ($-2$~\%). La prueba de que la competitividad ha funcionado son las exportaciones en volumen, un 56~\% por encima de 2007 (gráfico~8). El riesgo es que el diferencial de inflación y de costes de 2025-2026 se prolongue: \uline{en una unión monetaria, cada punto de inflación por encima de los socios es competitividad que se pierde sin que nadie devalúe}.
\end{examen}

\begin{complemento}{4.1 · Por qué se contesta así: de dónde viene y si persiste}
\que{Qué se aplica.} Guía, III.1: IPC frente a deflactor (inflación importada); la subyacente «dice si un choque de oferta se está extendiendo»; el efecto base, $\pi^{12}_t-\pi^{12}_{t-1}\approx\pi^m_t-\pi^m_{t-12}$; el triángulo de \autor{Gordon} (demanda, costes, inercia) \parencite{Gordon2011Phillips}; directos, indirectos y segunda ronda; transitoria o persistente y qué política. III.2: la convención del tipo de cambio real. II.1: CLU y margen.

\begin{tablacaja}
\centering\scriptsize
\caption{El triángulo de \autor{Gordon} con los datos de España en 2025-2026}\label{tab:gordon}
\begin{tabular}{>{\raggedright\arraybackslash}p{2.6cm}>{\raggedright\arraybackslash}p{7.2cm}>{\raggedright\arraybackslash}p{4.2cm}}
\toprule
\rowcolor{tceeazulpalido}
\textbf{Origen} & \textbf{Indicio en el cuadernillo} & \textbf{Lectura} \\
\midrule
Costes (oferta) & Brent +83~\% en dólares y +86~\% en euros (dic. 2025-sep. 2026, G7); CLU +4,4~\% (C3) & El detonante; común con la zona en la energía \\
Demanda & Brecha +2,0~\% (C1); consumo privado +0,67~pp en el 2T de 2026 (G1); vivienda +12,7~\% (C4) & Propio de España: explica el diferencial \\
Inercia & Subyacente del 3,0 al 3,5-3,7~\% (G6); salarios +4,4~\% con la productividad estancada (C3) & El riesgo de persistencia \\
\bottomrule
\end{tabular}
\fuente{Cuadernillo (C: cuadro; G: gráfico). Clasificación: guía, III.1 \parencite{Gordon2011Phillips}.}
\end{tablacaja}

\que{¿Transitoria o persistente?} Con la regla de la guía (III.1): \emph{transitoria} si se concentra en la energía con la subyacente contenida y salarios moderados; \emph{persistente} si la subyacente sube, los salarios se aceleran y la brecha es positiva. España tiene rasgos de la segunda; la zona del euro, de la primera. Eso justifica que la respuesta sea más intensa en España que en la zona, y como la monetaria es común, por la vía fiscal y de rentas.

\que{Error típico.} Atribuir el diferencial de 2026 «al petróleo» sin mirar la subyacente; o decir que el BCE debe «subir más por España»: la política monetaria responde a la zona del euro.
\end{complemento}

\begin{complemento}{4.2 · Los cálculos: diferencial, margen y competitividad}
\que{Descomponer el diferencial de septiembre de 2026} (gráfico~6): general $5{,}0-3{,}8=1{,}2$; subyacente $3{,}5-2{,}5=1{,}0$; general menos subyacente, $5{,}0-3{,}5=1{,}5$ en España y $3{,}8-2{,}5=1{,}3$ en la zona. No son contribuciones exactas (la subyacente no pesa lo mismo en los 2 índices), pero bastan para ver que el diferencial es sobre todo subyacente. Medias de enero-septiembre: general 3,6 frente a 2,8; subyacente 3,4 frente a 2,4.

\que{¿Es alcanzable el 3,0~\% previsto para 2026?} La media anual es la de los 12 meses: hasta septiembre suman 32,5, y para que la media fuera 3,0 los 3 últimos meses tendrían que promediar $(36-32{,}5)/3\approx1{,}2~\%$, lejos del 5,0~\% de septiembre.

\que{Inflación importada en 2022}: IPCA $-$ deflactor del PIB $=8{,}3-4{,}7=3{,}6$~puntos; relación real de intercambio, deflactores de exportación e importación: $1{,}146/1{,}224-1\approx-6{,}4~\%$ (empeora) y en 2023 $1{,}015/0{,}955-1\approx+6{,}3~\%$ (se recupera).

\begin{figuracaja}
\caption{España: CLU nominales, deflactor del PIB y margen, 2019-2025}\label{fig:s04}
\grafico[0.76\linewidth]{s04-margen}
\fuente{Cuadros~3 y~5 del cuadernillo (Eurostat). Margen: $\gamma_{1+\mu}\approx\gamma_P-\gamma_{\mathrm{CLU}}$ (guía, II.1). Qué hay que ver: en 2021-2022 los márgenes se amplían (el deflactor sube más que los CLU: las empresas trasladan más que el coste laboral) y la participación de los beneficios sube del 37,8 al 40,3~\% del VAB (cuadro~4); en 2025 se estrechan ($-1{,}4$): no hay «inflación de beneficios» al llegar el choque de 2026.}
\end{figuracaja}

\begin{tablacaja}
\centering\scriptsize
\caption{Costes y precios acumulados: España frente a la zona del euro (\%)}\label{tab:clu-acum}
\begin{tabular}{>{\raggedright\arraybackslash}p{4.4cm}*{4}{>{\centering\arraybackslash}p{1.9cm}}}
\toprule
\rowcolor{tceeazulpalido}
 & \textbf{2007-2013} & \textbf{2013-2019} & \textbf{2019-2025} & \textbf{2007-2025} \\
\midrule
CLU nominales, España & 1,2 & 5,0 & 26,8 & 34,7 \\
CLU nominales, zona del euro & 12,1 & 6,5 & 23,0 & 46,8 \\
IPCA, España & 13,4 & 3,4 & 21,5 & 42,6 \\
IPCA, zona del euro & 12,4 & 5,4 & 22,9 & 45,6 \\
Tipo de cambio efectivo real, CLU (variación) & $-10{,}0$ & $-1{,}9$ & 2,5 & $-9{,}4$ \\
\bottomrule
\end{tabular}
\fuente{Cálculo propio con los cuadros~3, 5 y~6 (\texttt{calculos.txt}). Tipo de cambio: un descenso es una depreciación real (ganancia de competitividad). Umbral del procedimiento de desequilibrios: $\pm3~\%$ en 3 años en la zona del euro (guía, III.2); 2022-2025: $98{,}9/99{,}8-1=-0{,}9~\%$, sin alerta.}
\end{tablacaja}

\que{Alternativas de redacción.} «La desinflación española de 2023 fue más rápida porque su inflación era más energética; la de 2026 es más alta porque es más \emph{doméstica}». \que{Qué añadir con más tiempo.} La inflación desigual (los hogares de renta baja gastan más en energía y alimentos, guía III.1) y las ayudas fiscales: temporales, focalizadas y que mantengan el incentivo a ahorrar energía (guía, VI.2). \que{Qué recortar en el examen.} El efecto base, la relación real de intercambio y el detalle del tipo de cambio con IPC: bastan la subyacente, el CLU y los 3 tiempos de la competitividad.
\end{complemento}

% Pregunta 5. Competitividad no precio y política comercial (bloque III.3) · peso bajo
\section*{Pregunta 5. Competitividad no precio y aranceles}

\begin{enunciado}
\textbf{5.} Valore la \textbf{competitividad no precio} de la economía española. ¿Cree que la política arancelaria de la actual administración de Estados Unidos ha tenido algún impacto sobre la economía española? ¿Qué medidas de política comercial y de internacionalización recomendaría?

\smallskip
\textbf{Peso bajo} · \tiempo{13} · guía: III.3 (no precio y política comercial), B.1 (aranceles de un tercero; canales de un choque externo) · cuadro 6; gráficos 8 y 9.
\end{enunciado}

\begin{examen}
\textit{La competitividad no precio de España es sólida pero no brillante: mantiene su cuota en las exportaciones mundiales y gana peso en servicios y tecnología, aunque invierte poco en innovación. Los aranceles de Estados Unidos han reducido las ventas a ese mercado, pero su efecto directo sobre la economía es pequeño.}

\rotulo{a) La competitividad no precio}
La cuota de España en las exportaciones mundiales de bienes y servicios se mantiene en torno al 2~\% desde 2007 (cuadro~6), algo notable con el auge de los emergentes, y las exportaciones en volumen crecen un 56~\% con un tipo de cambio real casi estable frente a los socios (gráfico~8): ganan mercados por algo más que el precio. Destacan los \textbf{servicios}, cuya cuota se recupera del 1,7~\% en 2020 al 2,7~\% en 2025, y no solo por el turismo: el saldo de los servicios no turísticos pasa de cero en 2007 al 2,6~\% del PIB en 2025 (cuadro~7). Las exportaciones de alta tecnología duplican su peso (del 4,2 al 8,5~\% entre 2007 y 2022). El punto débil es el \textbf{gasto en I+D}, del 1,50~\% del PIB frente al 2,26~\% de la UE, junto con un tejido de empresas pequeñas, que limita la innovación y la internacionalización.

\rotulo{b) El impacto de los aranceles de Estados Unidos}
Las exportaciones de bienes a Estados Unidos (suma de 12 meses, gráfico~9) caen un 11~\% entre marzo de 2025 y mayo de 2026, mientras las dirigidas al mundo se mantienen: hay un \textbf{efecto directo}, pero acotado, porque Estados Unidos solo absorbe en torno al 4~\% de las exportaciones de bienes (unos 17.000 millones de euros, el 1~\% del PIB). La pérdida, unos 2.000 millones en un año, equivale a algo más de una décima del PIB. Desde junio las ventas repuntan, en parte por el efecto base frente a los meses más afectados de 2025. Más relevantes pueden ser los \textbf{efectos indirectos}: la menor demanda de los socios europeos que exportan a Estados Unidos, la desviación hacia Europa de productos que ya no entran en ese mercado y, sobre todo, la \emph{incertidumbre}, que retrasa la inversión.

\rotulo{c) Medidas}
La política comercial es competencia de la UE, de modo que las recomendaciones tienen 2 planos. En el europeo, negociar con Estados Unidos, usar la defensa comercial frente a la desviación de comercio y ampliar los acuerdos con otras regiones para \textbf{diversificar mercados}. En el nacional, \textbf{internacionalizar} más empresas (promoción, financiación y seguro de crédito a la exportación, apoyo a las pymes para que exporten de forma regular) y, en el fondo, más \textbf{I+D} y \textbf{tamaño empresarial}, que son lo que sostiene la competitividad no precio.
\end{examen}

\begin{complemento}{5.1 · Por qué se contesta así, con los cálculos}
\que{Qué se aplica.} Guía, III.3 (cuota de mercado que se mantiene sin ganar competitividad precio: competitividad no precio; recomendaciones) y B.1 (efectos de los aranceles de un tercero: directo, cadenas de valor, desviación de comercio, incertidumbre y tipo de cambio).

\begin{figuracaja}
\caption{Los canales de un choque externo (infografía de la guía)}\label{fig:canales-choque}
\grafico[0.62\linewidth]{canales-choque-externo}
\fuente{Guía del dictamen, B.1. Qué hay que ver: para los aranceles de Estados Unidos se activan sobre todo el canal comercial y el de confianza; el de materias primas es el del petróleo (pregunta~4).}
\end{figuracaja}

\que{Cálculos} (gráfico~9 y \texttt{calculos.txt}): suma de 12 meses de las exportaciones a Estados Unidos, 18.660 millones en marzo de 2025 y 16.641 en mayo de 2026 ($-10{,}8~\%$); 17.412 en agosto de 2026 ($-6{,}7~\%$ frente a marzo de 2025). Al mundo, +0,6~\% en el mismo periodo. Pérdida: unos 2.020 millones, el $0{,}12~\%$ del PIB de 2025 (1.690 miles de millones). Es una cota del efecto \emph{bruto} en valor: el valor añadido perdido es menor (las exportaciones llevan contenido importado), pero los efectos indirectos lo amplían. Peso de Estados Unidos en las exportaciones de bienes (sumas de 12 meses): 4,7~\% a principios de 2025 y 4,2~\% en mayo de 2026.

\que{Servicios no turísticos} (cuadro~7): servicios menos viajes, $2{,}9-3{,}0\approx0{,}0$ en 2007 y $6{,}8-4{,}3=2{,}6$~\% del PIB en 2025.

\que{Errores típicos.} Exagerar el impacto («los aranceles hunden las exportaciones»): el dato agregado no lo muestra. Leer la suma de 12 meses como si fuera el mes: tarda en reflejar los cambios (y el repunte interanual de julio y agosto de 2026, +20 y +39~\%, se compara con los peores meses de 2025). Recomendar aranceles de represalia nacionales: la política comercial es de la UE.

\que{Qué añadir con más tiempo.} El reparto de las exportaciones por destinos (el cuadernillo solo da Estados Unidos y el total), para medir la exposición indirecta a través de los socios europeos, y el tipo de cambio del dólar, que cambia la competitividad frente a Estados Unidos. \que{Qué recortar en el examen.} Nada: es la pregunta corta; 3 párrafos de 3 o 4 frases.
\end{complemento}

% Pregunta 6. Cuentas públicas y política fiscal (bloque V.3-V.4) · peso alto
\section*{Pregunta 6. Cuentas públicas y política fiscal}

\begin{enunciado}
\textbf{6.} Describa la situación de las \textbf{cuentas públicas} y los principales rasgos de la \textbf{política fiscal}, su relación con el ciclo y su combinación con la política monetaria. ¿Considera que existen riesgos apreciables para la sostenibilidad de la deuda pública? ¿Cómo condicionarán las reglas fiscales europeas vigentes la política fiscal de los próximos años?

\smallskip
\textbf{Peso alto} · \tiempo{15} · guía: V.3 (saldos y tono, reglas, dinámica de la deuda, inflación y cuentas públicas), V.4 (\emph{policy mix}) · cuadros 1 y 9; gráficos 13 y 14.
\end{enunciado}

\begin{examen}
\textit{Las cuentas públicas han mejorado mucho desde la pandemia, pero siguen siendo el desequilibrio macroeconómico más claro de España: un déficit estructural del 3~\% del PIB con la economía por encima de su potencial y una deuda en torno al 100~\% que ha bajado gracias al crecimiento nominal más que al esfuerzo fiscal. La política fiscal ha pasado de procíclica a contracíclica, pero con un ajuste tímido, y las nuevas reglas europeas y el envejecimiento obligan a continuarlo.}

\rotulo{a) La situación de las cuentas públicas}
El déficit baja del 9,9~\% del PIB en 2020 al 2,4~\% en 2025, y la Comisión prevé el 2,4 y el 2,0~\% en 2026-2027 (cuadro~9). Con unos intereses del 2,4~\% del PIB, el saldo primario está equilibrado. La deuda, que pasó del 35,7~\% en 2007 a un máximo del 119,3~\% en 2020, está en el 100,5~\%. Pero el saldo total favorece la imagen: con la brecha en el +2,0~\%, el ciclo aporta en torno a 1,2 puntos, y el saldo ajustado por el ciclo es del $-3{,}6~\%$; sin medidas puntuales, el \textbf{saldo estructural} es del $-3{,}2~\%$ del PIB potencial, y el estructural primario, del $-0{,}8~\%$. Es decir, \uline{con la economía por encima de su potencial, las cuentas públicas aún tienen un déficit estructural del 3~\%}, por encima del 1,5~\% que fijan como referencia las reglas europeas. Los ingresos (42,9~\% del PIB) han subido casi 4 puntos desde 2019 (39,0~\%), y los gastos (45,3~\%) siguen 3 puntos por encima de los de antes de la pandemia (42,0~\%): parte del gasto extraordinario de 2020-2022 se ha quedado como gasto estructural.

La \textbf{inflación} ha ayudado a corto plazo, como predice la teoría: las bases tributarias son nominales y crecen con los precios (y, si los tramos del impuesto sobre la renta no se actualizan, la recaudación crece más que la renta), mientras los intereses solo suben a medida que se renueva la deuda (el tipo implícito pasa del 1,9 al 2,5~\% entre 2021 y 2025, pese a que el bono a 10 años llegó al 3,5~\% en 2023) y las pensiones y los salarios públicos se actualizan con retraso. A medio plazo el efecto se invierte: el gasto indexado y los intereses alcanzan a los ingresos.

\rotulo{b) Los rasgos de la política fiscal y su relación con el ciclo}
El tono se mide con la variación del saldo estructural primario, $\Delta b^{pe}$ (gráfico~13):
\begin{itemize}
\item en \textbf{2012-2013} la consolidación fue muy intensa (+2,9 y +2,5 puntos) con la brecha en su mínimo ($-8{,}7$ y $-9{,}6~\%$): una política \textbf{procíclica}, forzada por la pérdida de acceso a la financiación (diferencial de 435~pb);
\item en \textbf{2015-2019} fue expansiva, primero con la brecha negativa (contracíclica) y después con la brecha positiva ($-1{,}3$ puntos en 2019 con una brecha del +3,1~\%): procíclica de nuevo, porque no se aprovechó la expansión para reducir el déficit estructural y la economía llegó a la pandemia sin colchón;
\item en \textbf{2020-2022}, expansiva frente a la pandemia y a la crisis energética;
\item desde \textbf{2023}, contractiva con la brecha positiva (+1,3, +0,4 y +0,3 puntos en 2023-2025, y +0,4 en 2026-2027): por fin \textbf{contracíclica}, aunque \emph{tímida}, porque un déficit estructural del 3~\% que se reduce 3 o 4 décimas al año sigue siendo, por su nivel, una política holgada.
\end{itemize}
Los estabilizadores automáticos funcionan: en 2025 el ciclo mejora el saldo en 1,2 puntos, y en 2020 lo empeoró en 5.

\rotulo{c) La combinación con la política monetaria}
En 2026 la política monetaria pasa de algo expansiva a neutral para la zona del euro, pero sigue siendo expansiva para España (tipo real negativo, pregunta~4), y la fiscal es ligeramente contractiva. Las 2 van en la dirección correcta frente a una inflación alta y una brecha positiva: el \emph{policy mix} es \textbf{coherente}. Pero como la política monetaria es común y en España es más laxa que en la zona, \uline{la política fiscal es el único instrumento nacional para enfriar la demanda}, y debería apretar algo más de lo previsto; las ayudas frente a la energía, si las hay, deberían ser temporales y focalizadas en los hogares vulnerables. No siempre fue así: en 2018-2019 se combinó una política monetaria muy expansiva con una fiscal expansiva y una brecha positiva, y en 2022 las 2 seguían siendo expansivas con una inflación del 8~\%.

\rotulo{d) ¿Hay riesgos para la sostenibilidad de la deuda?}
La dinámica de la deuda es
\[
d_t-d_{t-1}=(g_t-\tau_t)+\frac{i_t-\gamma^N_t}{1+\gamma^N_t}\,d_{t-1}+\mathrm{ADD}_t .
\]
En 2025 el tipo implícito es del 2,5~\% y el crecimiento nominal, del 5,7~\% (gráfico~14): con $i-\gamma^N=-3{,}2$ puntos, el saldo primario que estabiliza la deuda es un \emph{déficit} del 3,1~\% del PIB. Con el saldo primario equilibrado la deuda debería haber bajado unos 3 puntos; bajó 0,9, lo que deja un ajuste entre déficit y deuda de unos 2 puntos (operaciones financieras que el cuadernillo no permite identificar). Entre 2020 y 2025 la deuda baja 18,8 puntos, y lo hace por el \textbf{efecto bola de nieve}, que la reduce en 31 puntos, frente a unos déficits primarios que suman 9: \uline{la deuda la ha bajado el crecimiento nominal, con la inflación incluida, no el ajuste fiscal}.

Por eso hay \textbf{riesgos apreciables a medio plazo}, aunque no inmediatos. El crecimiento nominal se moderará (5,1 y 4,5~\% previstos) y el tipo de la deuda nueva (el bono a 10 años está en el 3,9~\% en septiembre de 2026) ya supera al implícito (2,5~\%), de modo que el coste medio subirá a medida que se renueve la deuda. Si $i$ y $\gamma^N$ convergen hacia el 3,5~\%, harán falta superávits primarios para que la deuda baje. A ello se suma el gasto en pensiones y sanidad del envejecimiento (la tasa de dependencia sube del 24 al 31~\%). Lo mitigan que la deuda está en euros, un diferencial de solo 45~pb y una economía que crece.

\rotulo{e) Las reglas fiscales europeas}
Las reglas reformadas en 2024 mantienen las referencias del 3~\% de déficit y del 60~\% de deuda, pero se aplican con un solo indicador, una \textbf{senda de gasto primario neto} fijada en un plan fiscal-estructural de medio plazo (4 años, ampliables a 7 con reformas e inversiones), y con 2 salvaguardas: con una deuda superior al 90~\% del PIB, esta debe bajar al menos 1 punto al año de media, y el déficit estructural debe converger hacia un margen del 1,5~\%. Para España suponen 2 cosas: la deuda prevista (99,6 y 98,9~\%) baja 0,9 y 0,7 puntos al año, algo menos de lo que pide la salvaguarda, y el déficit estructural tiene que reducirse en torno a 1,7 puntos. Las reglas obligan, por tanto, a una \textbf{consolidación gradual y sostenida}, de algo más de medio punto al año, mejor por el lado del gasto corriente y ampliando bases que recortando inversión, con cláusulas de escape para circunstancias excepcionales.
\end{examen}

\begin{complemento}{6.1 · Los saldos, del total al estructural primario}
\que{Qué se aplica.} Guía, V.3: las 5 variables del saldo, $b^{pe}=b-\varepsilon\cdot\mathrm{OG}-\text{puntuales}+\text{intereses}$ y $\text{esfuerzo}=\Delta b^{pe}$; «cuidado con el signo»; con déficits muy altos, mirar también el nivel \parencite{EFB2021FiscalStance}; los estabilizadores, $b-b^{ca}=\varepsilon\cdot\mathrm{OG}$.

\begin{tablacaja}
\centering\scriptsize\setlength{\tabcolsep}{3.5pt}
\caption{Del saldo total al esfuerzo fiscal (\% del PIB o del PIB potencial)}\label{tab:saldos}
\begin{tabular}{>{\raggedright\arraybackslash}p{3.5cm}*{9}{>{\centering\arraybackslash}p{0.93cm}}}
\toprule
\rowcolor{tceeazulpalido}
 & \textbf{2019} & \textbf{2020} & \textbf{2021} & \textbf{2022} & \textbf{2023} & \textbf{2024} & \textbf{2025} & \cellcolor{tceegrisclaro}\textbf{2026\,P} & \cellcolor{tceegrisclaro}\textbf{2027\,P} \\
\midrule
Saldo total, $b$ & $-3{,}1$ & $-9{,}9$ & $-6{,}7$ & $-4{,}6$ & $-3{,}3$ & $-3{,}2$ & $-2{,}4$ & $-2{,}4$ & $-2{,}0$ \\
Brecha, OG & 3,1 & $-8{,}5$ & $-3{,}3$ & 0,9 & 1,0 & 1,9 & 2,0 & 1,7 & 1,2 \\
Ajustado por el ciclo, $b^{ca}$ & $-4{,}9$ & $-4{,}8$ & $-4{,}7$ & $-5{,}1$ & $-3{,}9$ & $-4{,}3$ & $-3{,}6$ & $-3{,}5$ & $-2{,}7$ \\
\hspace*{0.8em}$\varepsilon$ implícita, $(b-b^{ca})/\mathrm{OG}$ & 0,59 & 0,60 & 0,61 & 0,58 & 0,64 & 0,61 & 0,59 & 0,60 & 0,60 \\
Estructural, $b^{e}$ & $-4{,}5$ & $-4{,}3$ & $-4{,}6$ & $-5{,}1$ & $-3{,}8$ & $-3{,}5$ & $-3{,}2$ & $-2{,}9$ & $-2{,}5$ \\
Estructural primario, $b^{pe}$ & $-2{,}2$ & $-2{,}1$ & $-2{,}5$ & $-2{,}8$ & $-1{,}5$ & $-1{,}0$ & $-0{,}8$ & $-0{,}4$ & 0,0 \\
\rowcolor{tceesalvia}
Esfuerzo, $\Delta b^{pe}$ & $-1{,}3$ & 0,1 & $-0{,}4$ & $-0{,}3$ & 1,3 & 0,4 & 0,3 & 0,4 & 0,4 \\
Tono y ciclo & \tiny exp. procícl. & \tiny neutral* & \tiny exp. contrac. & \tiny exp. procícl. & \tiny contr. contrac. & \tiny contr. contrac. & \tiny contr. contrac. & \tiny contr. contrac. & \tiny contr. contrac. \\
\bottomrule
\end{tabular}
\fuente{Cuadros~1 y~9 del cuadernillo (AMECO, primavera de 2026). Exp.: expansiva; contr.: contractiva; procícl./contrac.: procíclica o contracíclica según el signo de la brecha. (*) En 2020 las medidas de la pandemia se tratan como temporales y no entran en el saldo estructural (guía, V.3, nota): el saldo estructural no recoge el gran estímulo de ese año. Qué hay que ver: $\varepsilon\approx0{,}6$ en todos los años, la semielasticidad de la Comisión; la diferencia entre $b^{ca}$ y $b^{e}$ (0,4 puntos en 2025) son medidas puntuales que aumentan el déficit; $b^{pe}\approx b^e+\text{intereses}$ ($-3{,}2+2{,}4=-0{,}8$).}
\end{tablacaja}

\que{Error típico.} Juzgar el tono por el saldo total (que mejora 0,8 puntos en 2025 sobre todo por el ciclo y por los intereses) o leer al revés el signo: un $b^{pe}$ que pasa de $-1{,}0$ a $-0{,}8$ \emph{sube}: es contractivo (guía, V.3, subrayado).
\end{complemento}

\begin{complemento}{6.2 · La dinámica de la deuda, con números}
\que{Qué se aplica.} Guía, V.3 (ecuación de la dinámica de la deuda, saldo primario estabilizador, «la diferencia entre el tipo y el crecimiento nominal decide») y «La inflación y las cuentas públicas» (a corto plazo mejora la ratio; a medio plazo el efecto se invierte) \parencite{Blanchard2019Deuda}.

\que{Alternativa de redacción} de la (d) en 2 frases: «La deuda ha bajado 19 puntos desde 2020 sobre todo porque el PIB nominal ha crecido mucho más que el tipo de interés, no por el ajuste. Con el tipo marginal (3,9~\%) ya por encima del implícito (2,5~\%) y el crecimiento nominal moderándose, ese viento de cola se agota: sin superávits primarios la deuda dejará de bajar».

\begin{figuracaja}
\caption{España: variación de la deuda pública y sus componentes, 2021-2027}\label{fig:s08}
\grafico[0.72\linewidth]{s08-deuda}
\fuente{Cuadro~9 del cuadernillo; previsiones de 2026-2027 de la Comisión (tipo implícito previsto: intereses entre la deuda del año anterior). El ajuste déficit-deuda es el residuo. Qué hay que ver: la bola de nieve, muy negativa con la inflación de 2021-2023, pierde fuerza año a año; en 2026-2027 la deuda solo bajaría 0,7-0,9 puntos, menos que el punto anual de la salvaguarda.}
\end{figuracaja}

\begin{tablacaja}
\centering\scriptsize
\caption{Por qué bajó la deuda: descomposición de 2021-2025 (puntos del PIB)}\label{tab:deuda}
\begin{tabular}{>{\raggedright\arraybackslash}p{4.6cm}*{6}{>{\centering\arraybackslash}p{1.25cm}}}
\toprule
\rowcolor{tceeazulpalido}
 & \textbf{2021} & \textbf{2022} & \textbf{2023} & \textbf{2024} & \textbf{2025} & \textbf{2021-2025} \\
\midrule
Tipo implícito, $i$ (\%) & 1,9 & 2,2 & 2,4 & 2,5 & 2,5 & \\
Crecimiento nominal, $\gamma^N$ (\%) & 9,4 & 11,4 & 8,8 & 6,8 & 5,7 & \\
Déficit primario & 4,6 & 2,3 & 0,9 & 0,8 & 0,0 & 8,6 \\
Bola de nieve, $\frac{i-\gamma^N}{1+\gamma^N}d_{t-1}$ & $-8{,}1$ & $-9{,}5$ & $-6{,}4$ & $-4{,}3$ & $-3{,}1$ & $-31{,}5$ \\
Ajuste déficit-deuda (residuo) & $-0{,}1$ & 0,8 & 1,5 & $-0{,}4$ & 2,2 & 4,1 \\
\rowcolor{tceesalvia}
Variación de la deuda, $\Delta d$ & $-3{,}6$ & $-6{,}4$ & $-4{,}0$ & $-3{,}9$ & $-0{,}9$ & $-18{,}8$ \\
\bottomrule
\end{tabular}
\fuente{Cálculo propio con el cuadro~9 (\texttt{calculos.txt}). Ejemplo, 2025: $\frac{0{,}0249-0{,}0572}{1{,}0572}\times101{,}4=-3{,}1$; $\mathrm{ADD}=-0{,}9-0{,}0-(-3{,}1)=+2{,}2$.}
\end{tablacaja}

\begin{tablacaja}
\centering\scriptsize
\caption{Saldo primario que estabiliza una deuda del 100~\% del PIB según $i$ y $\gamma^N$}\label{tab:escenarios}
\begin{tabular}{>{\raggedright\arraybackslash}p{6.0cm}*{3}{>{\centering\arraybackslash}p{1.6cm}}}
\toprule
\rowcolor{tceeazulpalido}
\textbf{Escenario} & $\boldsymbol{i}$ & $\boldsymbol{\gamma^N}$ & \textbf{Saldo primario} \\
\midrule
2025 (observado) & 2,5 & 5,7 & $-3{,}0$ \\
2027 (crecimiento previsto, mismo tipo) & 2,5 & 4,5 & $-1{,}9$ \\
El tipo sube hacia el del bono; el crecimiento se modera & 3,5 & 4,0 & $-0{,}5$ \\
$i=\gamma^N$ & 3,5 & 3,5 & 0,0 \\
Escenario adverso & 4,0 & 3,0 & +1,0 \\
\bottomrule
\end{tabular}
\fuente{Cálculo propio: $\tau-g=\frac{i-\gamma^N}{1+\gamma^N}\,d$ con $d=100$ (guía, A.3.5). Qué hay que ver: el margen que da hoy $i<\gamma^N$ se va estrechando; para bajar la deuda 1 punto al año en el escenario central hará falta un superávit primario de algo más de medio punto.}
\end{tablacaja}

\end{complemento}

\begin{complemento}{6.3 · \emph{Policy mix} y reglas fiscales: cómo se razona}
\begin{figuracaja}
\caption{Los 4 cuadrantes del \emph{policy mix} (infografía de la guía)}\label{fig:policy-mix}
\grafico[0.46\linewidth]{policy-mix}
\fuente{Guía del dictamen, V.4. España en 2026: monetaria expansiva para España aunque neutral para la zona, fiscal ligeramente contractiva; con inflación y brecha positiva, el cuadrante deseable es el de «ambas contractivas».}
\end{figuracaja}

\que{Qué se aplica.} Guía, V.4: «el \emph{policy mix} es coherente si cada política contribuye a los objetivos de la otra, no si las 2 tienen el mismo signo»; en una unión monetaria «el país solo controla la fiscal, que tiene que compensar si su ciclo va por delante del de la zona». V.3: las reglas de 2024 \parencite{RegUE2024Preventivo,RegUE2024Correctivo} y la cláusula para el gasto en defensa \parencite{CE2025Defensa}; recomendaciones por desequilibrio (VI.2).

\que{Las salvaguardas, con números.} Deuda: 100,5 $\to$ 99,6 $\to$ 98,9, es decir, $-0{,}9$ y $-0{,}7$ puntos; la salvaguarda pide $-1{,}0$ de media en el periodo de ajuste. Déficit estructural: de $-3{,}2$ a $-1{,}5$ son 1,7 puntos; a 0,4 puntos al año (el esfuerzo previsto), unos 4 años. Como referencia externa, el FMI proponía un ajuste de en torno a 0,5 puntos del PIB al año entre 2026 y 2030 \parencite{FMI2026ArticuloIV}.

\que{Errores típicos.} Decir que España «incumple» las reglas por tener una deuda del 100~\%: el valor de referencia del 60~\% no se exige de golpe, sino una senda. Y recomendar una consolidación brusca: con el choque energético frenando la demanda, mejor gradual y creíble, anunciada pronto (guía, V.3).

\que{Qué añadir con más tiempo.} Las necesidades brutas de financiación, la vida media de la deuda y quién la tiene (no vienen en el cuadernillo; guía V.3); la calidad del ajuste (proteger la inversión pública y la educación; ampliar bases antes que subir tipos). \que{Qué recortar en el examen.} La lista de etapas del tono (basta: procíclica en 2012-2013 y 2019, contracíclica y tímida desde 2023) y la historia del \emph{policy mix} anterior a 2026.
\end{complemento}

\begin{complemento}{6.4 · Contexto y fuentes: lo que el cuadernillo no trae}
\que{Las reglas, en sus textos.} El brazo preventivo (planes fiscales-estructurales de medio plazo y senda de gasto primario neto) está en el Reglamento (UE) 2024/1263 \parencite{RegUE2024Preventivo} y el correctivo (procedimiento de déficit excesivo) en el Reglamento (UE) 2024/1264 \parencite{RegUE2024Correctivo}; los valores del 3 y del 60~\% siguen en el Protocolo n.º 12 del Tratado \parencite{TFUE2012Protocolo12}.

\que{Defensa.} En 2025 la Comisión invitó a activar de forma coordinada la cláusula de escape nacional para el gasto en defensa, con un margen de hasta el 1,5~\% del PIB \parencite{CE2025Defensa} (guía, V.3). Es un riesgo de gasto y, a la vez, un margen dentro de las reglas: por eso va aquí y no en el texto de examen, que no tiene datos de gasto en defensa.

\que{El fin de los fondos europeos.} El Mecanismo de Recuperación y Resiliencia (Reglamento (UE) 2021/241) acaba en 2026: la inversión pública que financiaba tendrá que pagarse con recursos propios o reducirse, una presión adicional sobre el gasto a partir de 2027.

\que{La recomendación del FMI.} Un ajuste de en torno a 0,5 puntos del PIB al año entre 2026 y 2030 \parencite{FMI2026ArticuloIV} (guía, C.2), coherente con lo que da el cálculo del texto (algo más de medio punto al año).

\que{Lo que el cuadernillo no da y conviene decir que falta}: la vida media de la deuda, quién la tiene (residentes, no residentes, el Eurosistema) y las necesidades brutas de financiación (guía, V.3); el Tesoro los publica en sus informes, pero no en una serie descargable (ver \texttt{datos/fuentes-no-encontrado.md}).
\end{complemento}

% Pregunta 7. Desequilibrios, riesgos y recomendaciones (bloque VI) · peso alto
\clearpage
\section*{Pregunta 7. Desequilibrios, riesgos y recomendaciones}

\begin{enunciado}
\textbf{7.} Identifique los principales \textbf{desequilibrios y vulnerabilidades} de la economía española y los riesgos al alza y a la baja que rodean las previsiones para 2026 y 2027. Teniendo en cuenta las tendencias demográficas, ¿qué \textbf{reformas estructurales y políticas de oferta} recomendaría para sostener el crecimiento a medio y largo plazo?

\smallskip
\textbf{Peso alto} · \tiempo{15} · guía: VI.1 (fortalezas, desequilibrios y medidas; naturaleza, origen y gravedad), VI.2 (recomendaciones por desequilibrio), V.5 (políticas de oferta), D.8 (matriz), B.1 (demografía) · todo el cuadernillo.
\end{enunciado}

\begin{examen}
\textit{España afronta el final de la década con fortalezas que no tenía en 2007 (superávit exterior, sector privado desapalancado, banca rentable, empleo en máximos) y con 3 desequilibrios de fondo, por este orden: una productividad estancada, unas cuentas públicas con déficit estructural y deuda alta, y un paro estructural elevado; a ellos se suma, en 2026, un choque de precios que puede volverse persistente. El envejecimiento hace más urgente corregirlos, porque el crecimiento no podrá apoyarse indefinidamente en sumar trabajadores.}

\rotulo{a) Desequilibrios y vulnerabilidades}
Las \textbf{fortalezas} son claras: un crecimiento superior al de la zona del euro, la tasa de empleo más alta de la serie, un superávit exterior que se mantiene desde 2012, la deuda privada en su mínimo (hogares, 42,9~\%; empresas, 62,0~\% del PIB), una banca rentable y una prima de riesgo de 45~pb. Los desequilibrios, ordenados por su importancia:

\textbf{1.}~\textbf{Productividad estancada}: el PIB por hora crece un 0,33~\% anual desde 2019, el gasto en I+D es dos tercios del europeo y la inversión en maquinaria apenas crece. Es estructural y persistente, y es el más importante porque \uline{solo la productividad eleva la renta por habitante}, que desde 2007 crece menos que en la zona del euro.


\textbf{2.}~\textbf{Finanzas públicas}: un déficit estructural del 2,3~\% con la economía en su potencial, que las previsiones no reducen, y una deuda del 100,5~\%, que ha bajado gracias al crecimiento nominal. Es un problema de \emph{stock}, persistente, que se agravará con el envejecimiento y con la normalización del coste de la deuda.


\textbf{3.}~\textbf{Paro estructural}: el 10,5~\%, el doble que en la zona del euro, con un paro juvenil del 24,9~\% y muy pocas vacantes: un problema de emparejamiento y de cualificación más que de demanda.


\textbf{4.}~\textbf{Inflación y costes}: coyuntural, pero con riesgo de enquistarse; la inflación del 5,0~\% y unos costes laborales que crecen más que en la zona erosionan la competitividad ganada.

Los 2 primeros son problemas de \emph{stock} y de oferta, que no se corrigen solos con el ciclo; el tercero es en parte la histéresis de las crisis pasadas; el cuarto es de flujo, pero los problemas de flujo que persisten acaban siendo de \emph{stock}, y por eso conviene atajarlo pronto. Y 3 \textbf{vulnerabilidades}: la dependencia energética (el déficit de bienes llegó al 4,4~\% del PIB en 2022 y el superávit corriente se reducirá en 2026), la vivienda (precios a 2 dígitos con una oferta que no responde) y una posición exterior todavía por debajo del umbral europeo, aunque cada vez menos preocupante.

\rotulo{b) Riesgos de las previsiones}
Las previsiones del FMI (PIB del 2,5 y el 1,8~\%; inflación del 3,7 y el 2,7~\%) ya incorporan el choque del petróleo. Los riesgos para el \textbf{crecimiento} son, en 2026, equilibrados: el ritmo del primer semestre (0,6-0,7~\% trimestral) es coherente con la previsión, pero el petróleo, por encima de los 110 dólares en septiembre, resta renta real a los hogares y márgenes a las empresas, y empeora la balanza de bienes (en el choque de 2022 el déficit de bienes pasó del 1,7 al 4,4~\% del PIB). En 2027 dominan los riesgos \textbf{a la baja}: el encarecimiento del crédito (el Euríbor ha subido casi 1 punto y los tipos de los préstamos se reprecian deprisa), la debilidad de los socios por los aranceles, la consolidación fiscal pendiente y un posible ajuste de la vivienda. \textbf{Al alza}: la llegada de población, el turismo, un ahorro de los hogares todavía alto (11,9~\%) que da margen al consumo, y la inversión en intangibles. Para la \textbf{inflación}, los riesgos son al alza: la previsión anual (3,7~\%) exige que el último trimestre baje al 4~\% desde el 5,0~\% de septiembre, y la subyacente no cede.

\rotulo{c) Reformas estructurales y políticas de oferta}
\textbf{A corto plazo}, con una inflación del 5,0~\% y una subyacente del 3,5~\% (gráfico~6), la prioridad sería que el choque energético no se convierta en inflación persistente: una política fiscal que no añada demanda (ayudas, si las hay, temporales y focalizadas en los hogares vulnerables y que mantengan el incentivo a ahorrar energía), un pacto de rentas y una política macroprudencial preparada para frenar el crédito hipotecario si se acelera.\footnote{El FMI recomienda retirar las medidas energéticas al expirar y, en un escenario grave, sustituirlas por ayudas focalizadas que no alteren los precios \fmi{17}.} La política monetaria es común y no puede ajustarse a España.

\textbf{A medio y largo plazo}, el punto de partida es la \textbf{demografía}: la tasa de dependencia de los mayores sube del 24,0 al 31,2~\% entre 2007 y 2025 (cuadro~3), y el peso de la población en edad de trabajar ya no suma al crecimiento (cuadro~2~bis), que se ha apoyado en la inmigración. Para sostenerlo, haría falta que creciera la productividad y que trabajara más gente. Las recomendaciones, ordenadas por la importancia del desequilibrio que corrigen, con su objetivo y su riesgo:
\begin{itemize}
\item \textbf{Productividad}: con un PIB por hora que crece un 0,33~\% anual desde 2019 (cuadro~2~bis) y un gasto en I+D del 1,50~\% del PIB, frente al 2,26~\% de la UE (cuadro~6), convendría reforzar el capital humano (menos abandono escolar, formación profesional dual, recualificación), la I+D (incentivos fiscales, colaboración entre universidades y empresas, capital riesgo) y un entorno que permita a las empresas crecer (menos trabas regulatorias ligadas al tamaño, unidad de mercado, competencia en los servicios). Objetivo: un crecimiento más intensivo. Riesgo: sus efectos tardan años.\footnote{El FMI recomienda una reforma de la universidad que reduzca el desajuste de cualificaciones, suprimir umbrales fiscales y regulatorios ligados al tamaño, avanzar en la unidad de mercado («Régimen 20») y en la integración de los mercados de capitales europeos, y rediseñar el crédito fiscal a la I+D para que llegue a las empresas jóvenes \fmi{30-31}.}
\item \textbf{Mercado de trabajo}: con un paro del 10,5~\%, un paro juvenil del 24,9~\% y una tasa de vacantes de solo el 0,9~\% (cuadro~3, gráfico~3), serían útiles unas políticas activas de empleo focalizadas en jóvenes y parados de larga duración y evaluadas, una mejor intermediación de los servicios públicos de empleo, consolidar la menor temporalidad (15,4~\%) y salarios ligados a la productividad. Objetivo: bajar el paro estructural. Riesgo: la eficacia de las políticas activas depende de su diseño y su evaluación.\footnote{El FMI recomienda financiar los servicios públicos de empleo según su tasa de colocación, cotizaciones más altas para las empresas que despiden y contratan con más frecuencia, ampliar la compatibilidad del subsidio con el empleo y no subir el salario mínimo por encima del 60~\% del salario medio neto \fmi{28-29}.}
\item \textbf{Demografía y pensiones}: con una población extranjera que pasa del 10,3 al 14,1~\% del total desde 2019 y una tasa de dependencia al alza (cuadro~3), convendría integrar bien la inmigración (reconocimiento de cualificaciones, formación), elevar la actividad de los mayores y de las mujeres (jubilación flexible, servicios de cuidados) y asegurar la sostenibilidad de las pensiones ligándolas a la evolución demográfica. Objetivo: que la población activa siga creciendo aunque aumenten las jubilaciones.\footnote{El FMI propone, entre otras opciones, alargar el número de años cotizados que se usan para calcular la pensión \fmi{19}.}
\item \textbf{Consolidación fiscal gradual}: con un déficit estructural del 2,3~\% que las previsiones no reducen y una deuda del 100,5~\% (cuadro~9), unas décimas del PIB al año hasta llevar el déficit estructural por debajo del 1,5~\% y lograr superávits primarios, con revisiones del gasto, menos beneficios fiscales y bases más amplias, protegiendo la inversión pública. Riesgo: frenar la demanda si se hace deprisa.\footnote{El FMI recomienda 0,5 puntos al año entre 2026 y 2030, con medidas adicionales por un 1,5~\% del PIB, sobre todo armonizando los tipos reducidos del IVA \fmi{17-18}.}
\item \textbf{Vivienda}: los precios suben un 12,7~\% en 2025 (un 9,7~\% en términos reales), por encima del umbral del 9~\% del procedimiento de desequilibrios, pero, a diferencia de 2007, sin que los impulse el crédito: la deuda de los hogares está en mínimos de la serie (42,9~\% del PIB; cuadro~4, gráfico~4). Al mismo tiempo, la población y el empleo crecen con fuerza (un 0,7 y un 2,6~\% en 2025; cuadro~3) y la inversión en vivienda apenas aumenta (un 6,0~\% desde 2019, frente a un 13,4~\% del total de la inversión; cuadro~1). Todo apunta a un desajuste entre una demanda que crece y una oferta que, al menos a corto plazo, no la absorbe. Como el precio depende de ambos lados del mercado, convendría actuar sobre los 2: por el lado de la oferta, facilitando que se construya más y ampliando el parque de alquiler asequible; por el lado de la demanda, moderando la demanda no residencial o especulativa (por ejemplo, con la fiscalidad) y vigilando el crédito hipotecario con instrumentos macroprudenciales por si se acelera. Serían menos aconsejables las ayudas generales a la compra, que con una oferta poco flexible tenderían a trasladarse a los precios. La vivienda cara podría estar dificultando también la movilidad geográfica y el emparejamiento en el mercado de trabajo (las vacantes están en su máximo desde 2011 con un paro todavía del 10,5~\%; cuadro~3, gráfico~3).\footnote{El FMI recomienda acelerar el desarrollo urbanístico, simplificar las licencias, retomar la reforma de la ley del suelo y que el control de alquileres sea temporal y esté condicionado a aumentar la oferta \fmi{29, 81}.}
\item \textbf{Energía}: con un déficit de bienes que llegó al 4,4~\% del PIB en el choque de 2022 y que las previsiones sitúan de nuevo en el 4,3-4,5~\% en 2026-2027 (cuadro~7), convendría diversificar las fuentes, acelerar las renovables y reforzar las interconexiones y la eficiencia, para reducir la exposición a choques como los de 2022 y 2026.\footnote{El FMI señala que el peso creciente de las renovables amortigua el traslado del precio del gas a la electricidad \fmi{11}.}
\end{itemize}
En conclusión, España tiene hoy un margen que no tuvo en 2008 (superávit exterior y deuda privada baja; cuadros~4 y~7) para \uline{aprovechar la expansión y corregir los desequilibrios públicos y de productividad}; cómo lo haga decidirá si su crecimiento se traduce en renta por habitante.
\end{examen}

\begin{complemento}{7.1 · La matriz de fortalezas y debilidades, rellena}
\que{Qué se aplica.} Guía, D.8: la matriz se va rellenando durante el examen en una hoja aparte y sirve para 3 cosas: anticipar el diagnóstico en la pregunta~1, no contradecirse y contestar la última pregunta. VI.1: \emph{qué va bien, qué hay que cambiar (2 o 3 desequilibrios ordenados) y cómo}; las 3 preguntas del FMI (naturaleza, origen y gravedad) \parencite{FMI2013CaseStudy}.

\begin{tablacaja}
\centering\scriptsize
\renewcommand{\arraystretch}{0.98}\setlength{\tabcolsep}{4pt}
\caption{Matriz de fortalezas y debilidades de España (examen de ejemplo)}\label{tab:matriz-espana}
\begin{tabular}{>{\raggedright\arraybackslash}p{2.0cm}>{\raggedright\arraybackslash}p{3.3cm}>{\raggedright\arraybackslash}p{3.6cm}>{\centering\arraybackslash}p{0.7cm}>{\centering\arraybackslash}p{0.8cm}>{\raggedright\arraybackslash}p{3.4cm}}
\toprule
\rowcolor{tceeazulpalido}
\textbf{Bloque} & \textbf{\textcolor{tceeverde}{Fortaleza}} & \textbf{\textcolor{tceerojo}{Debilidad}} & \textbf{Peso} & \textbf{Tend.} & \textbf{Recomendación} \\
\midrule
Perfil cíclico & PIB +10,0~\% en 2019-2025 (zona: +6,8~\%) & Crecimiento extensivo: PIB por habitante +5,0~\% & ** & $\nearrow$ & Políticas de oferta (productividad) \\
Trabajo & Empleo 67,0~\%; temporalidad 15,4~\% & Paro 10,5~\%; juvenil 24,9~\%; PIB por hora +0,33~\%/año & *** & $\nearrow$ & Políticas activas, emparejamiento, capital humano \\
Capital y crédito & Deuda privada en mínimos; intangibles +47~\% & Maquinaria +2,7~\%; vivienda +12,7~\% & * & $\rightarrow$ & Vivienda: oferta y demanda; incentivos a invertir \\
Precios & Salida rápida de 2022-2023 & IPCA 5,0~\% (zona 3,8); CLU +4,4~\% & ** & $\searrow$ & Pacto de rentas; fiscal que no empuje \\
No precio & Cuota mundial estable; servicios & I+D 1,50~\% (UE 2,26) & * & $\rightarrow$ & I+D, tamaño, diversificar mercados \\
Exterior & Superávit desde 2012; NENDI $-9{,}7~\%$ & PIIN $-42{,}7~\%$; energía & * & $\nearrow$ & Política energética \\
Monetaria & Transmisión rápida; diferencial 45~pb & Tipo real más bajo en España & * & $\rightarrow$ & (común; compensa la fiscal) \\
Banca & ROE 14,2~\%; dudosos 2,4~\% & CET1 13,9~\% (UE 16,4) & * & $\nearrow$ & Colchones; límites LTV/LTI si hace falta \\
Fiscal & Saldo sin intereses equilibrado; deuda a la baja & Déficit estructural 2,3~\%; deuda 100,5~\% & *** & $\nearrow$ & Consolidación gradual con la senda de gasto \\
Demografía & Inmigración (extranjeros 14,1~\%) & Dependencia 31,2~\% & ** & $\searrow$ & Integración, actividad, pensiones \\
\bottomrule
\end{tabular}
\fuente{Elaboración propia con el cuadernillo; plantilla de la guía (D.8). Peso: importancia del desequilibrio (*** máxima). Tendencia: $\nearrow$ mejora, $\rightarrow$ estable, $\searrow$ empeora.}
\end{tablacaja}
\end{complemento}

\begin{complemento}{7.2 · Los umbrales oficiales: qué se enciende en España}

\begin{tablacaja}
\centering\scriptsize
\caption{España frente a los umbrales del procedimiento de desequilibrios y del Tratado}\label{tab:pdm}
\begin{tabular}{>{\raggedright\arraybackslash}p{4.4cm}>{\centering\arraybackslash}p{2.6cm}>{\centering\arraybackslash}p{3.2cm}>{\centering\arraybackslash}p{2.4cm}}
\toprule
\rowcolor{tceeazulpalido}
\textbf{Indicador} & \textbf{Umbral} & \textbf{España} & \textbf{¿Alerta?} \\
\midrule
Saldo corriente, media de 3 años & $-4$ / $+6~\%$ & $(2{,}7+3{,}3+3{,}2)/3=3{,}1$ & \textcolor{tceeverde}{no} \\
PIIN & $-35~\%$ & $-42{,}7~\%$ & \textcolor{tceerojo}{sí} (mejorando) \\
Tipo de cambio efectivo real, 3 años & $\pm3~\%$ & $-0{,}9~\%$ (2022-2025) & \textcolor{tceeverde}{no} \\
CLU nominales, 3 años & $+9~\%$ & $\approx+13{,}9~\%$ (2023-2025)* & \textcolor{tceerojo}{sí} \\
Precio de la vivienda (nominal) & $+9~\%$ & $+12{,}7~\%$ & \textcolor{tceerojo}{sí} \\
Deuda de hogares / de empresas & 55 / 85~\% & 42,9 / 62,0~\% & \textcolor{tceeverde}{no} \\
Tasa de paro & 10~\% & 10,5~\% & \textcolor{tceerojo}{sí} (por poco) \\
Saldo público (Tratado) & $-3~\%$ & $-2{,}4~\%$ & \textcolor{tceeverde}{no} \\
Deuda pública (Tratado) & 60~\% & 100,5~\% & \textcolor{tceerojo}{sí} \\
\bottomrule
\end{tabular}
\fuente{Umbrales: \textcite{CE2025AMR2026,TFUE2012Protocolo12} (guía, A.1); datos: cuadros 3, 4, 6, 7 y 9. (*) El umbral se refiere a los CLU por hora; con los CLU por persona del cuadernillo, $1{,}059\times1{,}030\times1{,}044-1\approx13{,}9~\%$, una aproximación. Qué hay que ver: las alertas están en los costes, la vivienda, el paro y la deuda pública: exactamente los desequilibrios de la respuesta. Son señales, no obligaciones (guía, A).}
\end{tablacaja}
\begin{figuracaja}
\caption{Los umbrales oficiales de un vistazo, con la posición de España}\label{fig:umbrales}
\grafico[0.66\linewidth]{s09-umbrales-espana}
\fuente{Infografía de la guía (A, umbrales del procedimiento de desequilibrios y del Tratado) con el último dato de España del cuadernillo. Qué hay que ver: los rombos en rojo (PIIN, paro, CLU, vivienda, deuda pública) son los desequilibrios de la respuesta; los demás, en verde.}
\end{figuracaja}
\end{complemento}

\begin{complemento}{7.3 · Riesgos, recomendaciones que responden a los datos y qué recortar}
\que{Del dato a la recomendación} (guía, VI.1: «puntúan las recomendaciones que responden a los datos del cuadro»): cada medida de la respuesta sale de un indicio (gráfico~\ref{fig:dato-recomendacion}). Ejemplo: \emph{PIB por hora +0,33~\% e I+D del 1,50~\%} $\to$ productividad estancada $\to$ I+D y tamaño empresarial $\to$ objetivo: crecimiento intensivo $\to$ riesgo: tarda.

\begin{figuracaja}
\caption{Del dato a la recomendación (infografía de la guía)}\label{fig:dato-recomendacion}
\grafico[0.56\linewidth]{esquema-dato-recomendacion}
\fuente{Guía del dictamen, VI.1. Qué hay que ver: siempre el mismo camino, del indicio al desequilibrio, a la medida, a su objetivo y a su riesgo.}
\end{figuracaja}

\begin{tablacaja}
\centering\scriptsize
\caption{Riesgos de las previsiones de 2026-2027 y su indicio en el cuadernillo}\label{tab:riesgos}
\begin{tabular}{>{\raggedright\arraybackslash}p{1.6cm}>{\raggedright\arraybackslash}p{5.2cm}>{\raggedright\arraybackslash}p{7.0cm}}
\toprule
\rowcolor{tceeazulpalido}
\textbf{Sentido} & \textbf{Riesgo} & \textbf{Indicio} \\
\midrule
A la baja & Petróleo caro y persistente & Brent de 62,5 a 114,2 dólares (G7); bienes $-4{,}4~\%$ del PIB en el choque de 2022 (C7) \\
 & Endurecimiento financiero & Depósito 2,50~\%; Euríbor +98~pb desde diciembre de 2025; hipotecas 3,04~\% (G12) \\
 & Aranceles y socios & Exportaciones a EE.~UU. $-11~\%$ (G9) \\
 & Menos transferencias de capital & Cuenta de capital: 1,3~\% del PIB en 2026 y 0,2~\% en 2027 (C7) \\
 & Vivienda & Precio +12,7~\% con poca inversión residencial (C1, C4) \\
Al alza & Población y turismo & Población +1,0~\% al año; viajes 4,3~\% del PIB (C3, C7) \\
 & Consumo & Ahorro de los hogares 11,9~\%, el doble que en 2007 (C4) \\
 & Ritmo de 2026 & 0,6-0,7~\% trimestral en el primer semestre (G1): $\approx$2,5~\% anual \\
Inflación (al alza) & Segunda ronda & Media de enero-septiembre 3,6~\% con septiembre en el 5,0; previsión 3,7~\% (C5); subyacente 3,5~\% (G6) \\
\bottomrule
\end{tabular}
\fuente{Cuadernillo (C: cuadro; G: gráfico).}
\end{tablacaja}

\que{Errores típicos} (guía, VI.1 y D.6): recomendaciones genéricas que valen para cualquier país («hacer reformas estructurales»); recetas tajantes o demasiado precisas; olvidar la demografía, que el enunciado pide expresamente; no ordenar los desequilibrios. \que{Alternativa de redacción} del cierre: «España tiene que aprovechar una expansión sin desequilibrios privados para corregir los públicos y elevar la productividad, antes de que el envejecimiento reduzca el crecimiento de la población activa». \que{Qué añadir con más tiempo.} El caso de España del FMI (guía, C.2): ajuste fiscal de unos 0,5 puntos al año entre 2026 y 2030, políticas activas y de vivienda, y reformas de productividad \parencite{FMI2026ArticuloIV}. \que{Qué recortar en el examen.} Las vulnerabilidades (1 frase) y 2 de las 6 recomendaciones (energía y vivienda, ya tratadas en las preguntas 1, 3, 4 y 5).
\end{complemento}

\begin{complemento}{7.4 · Contexto y fuentes: lo que el cuadernillo no trae}
\que{Riesgos que el cuadernillo no muestra.} El final del Mecanismo de Recuperación y Resiliencia en 2026 (Reglamento (UE) 2021/241), que ha financiado buena parte de la inversión en intangibles e infraestructuras, es un riesgo a la baja para 2027; el gasto en defensa (cláusula de escape de 2025, \cite{CE2025Defensa}), una presión sobre las cuentas públicas.

\que{El diagnóstico del FMI.} La consulta del artículo~IV de 2026 resume que España «siguió creciendo más que la zona del euro pese al impacto adverso de la guerra en Oriente Medio» y recomienda un ajuste fiscal gradual, políticas activas y de vivienda que mejoren el emparejamiento y reformas de productividad (capital humano, tamaño empresarial, I+D) \parencite{FMI2026ArticuloIV} (guía, C.2): las mismas líneas que salen de los datos en el texto de examen.

\que{Envejecimiento.} La Comisión proyecta un aumento del gasto ligado a la edad (pensiones, sanidad, cuidados) en la mayoría de los países de la UE \parencite{CE2024Envejecimiento}; con la tasa de dependencia de los mayores en el 31,2~\% y subiendo (cuadro~3), es el argumento de fondo para empezar ya la consolidación y las reformas de actividad.

\que{Las reformas ya hechas.} La reforma laboral de 2021 (Real Decreto-ley 32/2021, pregunta~2) es el ejemplo de una reforma cuyo efecto se ve en el cuadernillo (temporalidad del 25,4 al 15,4~\%): sirve para defender que las reformas de oferta funcionan, aunque sus efectos sobre la productividad tarden más en verse.
\end{complemento}

\begin{complemento}{7.5 · El diagnóstico del FMI (artículo~IV de 2026): riesgos y reformas}
\que{La matriz de riesgos} \fmi{52-56}. Con probabilidad alta: las tensiones geopolíticas (el conflicto de Oriente Medio es el principal riesgo a la baja), el proteccionismo, la volatilidad de las materias primas (España importa en torno al 70~\% de la energía que usa), una corrección brusca ligada a la inteligencia artificial, la subida de los tipos por la deuda pública mundial y los ciberataques. Riesgos internos: la fragmentación política (impacto alto), la ejecución de los fondos europeos (unos desembolsos del 1,2~\% del PIB en 2026) y el auge de la vivienda. Al alza: más turismo, más desahorro de los hogares, más inmigración (el potencial de 2030 sería 0,1 puntos mayor si se mantuviera la de 2025) y la inteligencia artificial (0,1-0,2 puntos más de productividad al año) \fmi{14}. \textbf{Coincide} con la tabla de riesgos de la solución y añade 2 que el cuadernillo no permite ver: la fragmentación política y la ejecución de los fondos.

\que{Las reformas} \fmi{28-31}. Mercado de trabajo (cuadro~2.5); vivienda: acelerar el desarrollo urbanístico, simplificar las licencias, recuperar la reforma de la ley del suelo y que el control de alquileres sea temporal y condicionado a aumentar la oferta (en Cataluña rebajó los alquileres un 4-6~\%, pero redujo los anuncios de alquiler entre un 15 y un 30~\%) \fmi{29, 81}; productividad: unidad de mercado («Régimen 20»), integración de los mercados de capitales europeos, reforma de la universidad para reducir el desajuste de cualificaciones, menos umbrales regulatorios ligados al tamaño y un crédito fiscal a la I+D más sencillo y dirigido a las empresas jóvenes, que podría elevar el crecimiento de la productividad total de los factores en un cuarto de punto \fmi{30-31}.

\textbf{Coincide} con las 6 líneas de la solución; \textbf{añade} detalle (el IVA y las pensiones en lo fiscal, el crédito fiscal a la I+D en la innovación) y una advertencia de economía política: con la fragmentación parlamentaria, el riesgo de no cumplir los planes es alto \fmi{34}.
\end{complemento}

\begin{complemento}{7.6 · La vivienda: ¿problema de oferta o de demanda? Las tijeras de \autor{Marshall}}
\que{Los 2 filos.} \autor{Marshall} comparó la oferta y la demanda con los 2 filos de unas tijeras: preguntarse cuál de los 2 determina el precio es como preguntarse qué filo corta el papel \parencite[libro~V, cap.~III]{Marshall1890}. Decir que la vivienda «es solo un problema de oferta» elige un filo sin pruebas; el cuadernillo muestra los 2: una demanda que crece con la población y el empleo y una inversión residencial que apenas responde (texto de examen, apartado~c).

\que{Qué filo conviene mover depende de las elasticidades.} A corto plazo el parque de viviendas es casi fijo (las viviendas terminadas están en un 14~\% del nivel de 2007 \fmi{47}). Con una oferta tan rígida:
\begin{itemize}
\item las medidas que \textbf{reducen la demanda} bajan la presión sobre el precio;
\item las que la \textbf{apoyan} sin ampliar la oferta se trasladan en buena parte al precio (se \emph{capitalizan}).
\end{itemize}
A largo plazo la oferta es más elástica, y ahí es donde rinden las medidas de oferta (suelo, licencias, seguridad jurídica), que tardan años.

\que{La evidencia sobre medidas de demanda}, con sus fuentes:
\begin{itemize}
\item \emph{Ayudas generales a la compra}: el programa británico \emph{Help to Buy} elevó los precios en Londres, donde la oferta está muy restringida, más que el valor de la subvención, sin efecto apreciable sobre la construcción \parencite{Carozzi2024HelpToBuy}.
\item \emph{Ayudas al alquiler}: en Francia, la ampliación de las ayudas a la vivienda de principios de los años 90 elevó los alquileres sin mejorar la calidad de las viviendas \parencite{Fack2006Ayudas}; según la reseña del Banco de Francia, la subida de los alquileres equivalió al 78~\% del importe de las ayudas \parencite{BdF2023AyudasVivienda}.
\item \emph{Fiscalidad de la vivienda vacía}: el impuesto francés sobre la vivienda vacía redujo su número en torno a un 13~\% entre 1997 y 2001, sobre todo la vacía de larga duración, y la mayoría pasó a ser vivienda principal \parencite{Segu2020Vacantes}.
\item \emph{Pisos turísticos}: en Barcelona, la actividad de Airbnb elevó los alquileres un 1,9~\% y los precios de compraventa un 4,6~\% en el barrio medio, y hasta un 7 y un 17~\% en los barrios con más actividad \parencite{GarciaLopez2020Airbnb}: limitarla podría aliviar la presión allí donde se concentra, aunque el efecto en el agregado sería menor.
\item \emph{Límites al crédito (préstamo sobre el valor o sobre la renta)}: frenan el crédito a los hogares de forma significativa; su efecto sobre los precios es más moderado \parencite{Alam2019Macroprudencial}. El FMI los recomienda para España, sobre todo para contener riesgos financieros \fmi{25-26}.
\item \emph{Compras de no residentes}: según el Banco de España, esta demanda contribuye a la subida de los precios y se concentra en Illes Balears, la Comunitat Valenciana y las Islas Canarias, donde supuso entre el 20 y el 25~\% de las compraventas en 2023 \parencite[pp.~246-247]{BdE2024Vivienda}. Una medida sobre ella actuaría, por tanto, sobre todo en esas zonas; su efecto en los precios dependería del diseño y no se cuantifica aquí.
\end{itemize}

\que{La localización.} El problema no es igual en todas partes: los problemas de acceso se concentran en las áreas urbanas y turísticas, y la creación de hogares, en las comunidades más pobladas y las grandes ciudades \parencite[pp.~232 y 242-243]{BdE2024Vivienda}. Si en unas falta vivienda y en otras hay vivienda vacía, las medidas eficaces serían locales (suelo y licencias donde falta; fiscalidad de la vacía donde abunda), y una media nacional escondería el problema. El FMI añade que la presión se ha extendido más allá de las grandes ciudades y la costa \fmi{8}.

\que{Error típico.} Recomendar «más ayudas a la compra para los jóvenes» con una oferta rígida: con los datos anteriores, buena parte acabaría en el precio. Y el contrario: descartar toda medida de demanda porque «el problema es la oferta».
\end{complemento}

% Pregunta 8. Sistema bancario, estabilidad financiera y macroprudencial (bloque V.2) · peso bajo
\section*{Pregunta 8. Sistema bancario y estabilidad financiera}

\begin{enunciado}
\textbf{8.} Valore la situación del \textbf{sistema bancario} español en términos de solvencia, calidad de los activos y rentabilidad. ¿Qué riesgos para la \textbf{estabilidad financiera} aprecia? ¿Qué medidas de política macroprudencial considera oportunas?

\smallskip
\textbf{Peso bajo} · \tiempo{13} · guía: V.2 (endeudamiento, indicadores bancarios, riesgos, macroprudencial; semáforo de la banca) · cuadros 4 y 8; gráfico 4.
\end{enunciado}

\begin{examen}
\textit{El sistema financiero español parece resiliente: la banca es muy rentable, tiene poca morosidad y cumple con holgura los requisitos de capital, aunque con menos capital de máxima calidad que la media europea, y la deuda privada está en mínimos. El riesgo que hay que vigilar es la vivienda.}

\rotulo{a) Solvencia, calidad de los activos y rentabilidad}
La \textbf{solvencia} es suficiente, pero no holgada: el capital ordinario de nivel 1 (CET1) es el 13,9~\% de los activos ponderados por riesgo y el capital total, el 18,1~\% (cuadro~8), por encima de los requisitos medios del Mecanismo Único de Supervisión (en torno al 11~\% de CET1), pero por debajo de la media de la UE (16,4~\%). La \textbf{calidad de los activos} es buena: los préstamos dudosos bajan del 8,0~\% en 2014 al 2,4~\%, por debajo del 3~\% que se considera bajo, aunque algo por encima de la media europea (2,0~\%). La \textbf{rentabilidad} es alta: un ROE del 14,2~\% (9,3~\% en la UE), por encima del coste del capital, lejos del $-25{,}6~\%$ de 2012. No se debe a un mayor apalancamiento, sino a un activo más rentable (ROA del 0,96~\% frente al 0,71~\%): el margen de intereses se beneficia de los tipos más altos y de un crédito a tipo variable que se reprecia deprisa.

\rotulo{b) Riesgos para la estabilidad financiera}
El principal es la \textbf{vivienda}: su precio sube un 12,7~\% en 2025, por encima del umbral del 9~\% del procedimiento de desequilibrios, y crece más que la inversión residencial (más precio que cantidad). Lo mitiga que no hay un auge del crédito: la deuda de los hogares está en el 42,9~\% del PIB y el crédito a hogares crece un 4,5~\%. El segundo es el \textbf{endurecimiento monetario} unido a una desaceleración: con mucho crédito a tipo variable, la subida de los tipos se traslada enseguida a las cuotas y podría elevar la morosidad, que va retrasada respecto al ciclo. Y, en las empresas, el encarecimiento de la energía y los aranceles. El bucle bancario-soberano preocupa poco con un diferencial de 45~pb.

\rotulo{c) Medidas macroprudenciales}
En una fase alcista con la vivienda al alza, lo prudente es \textbf{acumular capital ahora}: mantener activado el colchón de capital anticíclico y elevarlo si el crédito se acelera, y moderar el reparto de dividendos para acercar el CET1 a la media europea. Para la vivienda, tener preparados \textbf{límites a los prestatarios} (préstamo sobre el valor de la vivienda y sobre la renta) y activarlos si el crédito hipotecario se dispara. El problema de fondo de la vivienda, la oferta, no lo resuelve la política macroprudencial (pregunta~10).
\end{examen}

\begin{complemento}{8.1 · Por qué se contesta así: el semáforo y el orden de la guía}
\que{Qué se aplica.} Guía, V.2: empezar con una frase que resuma el estado del sistema y seguir el orden \emph{endeudamiento, banca, riesgos y políticas}; las ratios con sus referencias (gráfico~\ref{fig:semaforo}); $\mathrm{ROE}=\mathrm{ROA}\times\text{activo}/\text{patrimonio}$; la vivienda como «principal activo, principal deuda y garantía del crédito»; instrumentos macroprudenciales en la fase alcista.

\begin{figuracaja}
\caption{Semáforo de la banca: referencias de las ratios principales (infografía de la guía)}\label{fig:semaforo}
\grafico[0.58\linewidth]{esquema-semaforo-banca}
\fuente{Guía del dictamen, V.2. España en 2025 (cuadro~8): CET1 13,9~\% (\textcolor{tceedorado}{\textbf{vigilar}}, cerca de holgado); dudosos 2,4~\% (\textcolor{tceeverde}{\textbf{holgado}}); ROE 14,2~\% (\textcolor{tceeverde}{\textbf{holgado}}). La liquidez y la eficiencia no vienen en el cuadernillo.}
\end{figuracaja}

\que{Cálculos} (cuadro~8): apalancamiento implícito, $\mathrm{ROE}/\mathrm{ROA}=14{,}2/0{,}96\approx14{,}8$ (el patrimonio es el 6,8~\% del activo); en la UE, $9{,}3/0{,}71\approx13{,}2$. España apalanca algo más, pero la diferencia de ROE (4,9 puntos) se explica sobre todo por el ROA (0,96 frente a 0,71). Margen de CET1 sobre los requisitos medios del MUS (11,2~\%, guía V.2 \parencite{BCE2025SREP}): unos 2,7 puntos.

\que{Más información de la que trae el cuadernillo.} El colchón de capital anticíclico de España es del 1~\% desde octubre de 2026 \parencite{BdE2025CCA} (guía, tabla de Basilea~III), y las estadísticas supervisoras del Banco de España daban a mediados de 2026 un ROE del 16~\% y un CET1 del 14~\% \parencite{BdE2026EstadSupervisoras}. Con ese dato, la recomendación de (c) es \emph{mantenerlo y estar listos para subirlo}, no activarlo.

\que{Errores típicos.} Confundir el ROA bancario con el de una empresa no financiera; decir que la banca «no cumple» por tener menos CET1 que la media (cumple con margen); recomendar medidas de crisis (liberar colchones) en una fase alcista. \que{Qué añadir con más tiempo.} La brecha crédito-PIB (no viene; el flujo de crédito del 2,6~\% del PIB sugiere que es negativa, guía V.2) y la liquidez (LCR, créditos sobre depósitos). \que{Qué recortar en el examen.} Nada: es la otra pregunta corta.
\end{complemento}

% Pregunta 9. Cuentas públicas y política fiscal (bloque V.3-V.4) · peso alto
\section*{Pregunta 9. Cuentas públicas y política fiscal}

\begin{enunciado}
\textbf{9.} Describa la situación de las \textbf{cuentas públicas} y los principales rasgos de la \textbf{política fiscal}, su relación con el ciclo y su combinación con la política monetaria. ¿Considera que existen riesgos apreciables para la sostenibilidad de la deuda pública? ¿Cómo condicionarán las reglas fiscales europeas vigentes la política fiscal de los próximos años?

\smallskip
\textbf{Peso alto} · \tiempo{21} · guía: V.3 (saldos y tono, reglas, dinámica de la deuda, inflación y cuentas públicas), V.4 (\emph{policy mix}) · cuadros 1 y 9; gráficos 13 y 14.
\end{enunciado}

\begin{examen}
\textit{Las cuentas públicas han mejorado mucho desde la pandemia, pero siguen siendo el desequilibrio macroeconómico más claro de España: un déficit estructural del 3~\% del PIB con la economía por encima de su potencial y una deuda en torno al 100~\% que ha bajado gracias al crecimiento nominal más que al esfuerzo fiscal. La política fiscal ha pasado de procíclica a contracíclica, pero con un ajuste tímido, y las nuevas reglas europeas y el envejecimiento obligan a continuarlo.}

\rotulo{a) La situación de las cuentas públicas}
El déficit baja del 9,9~\% del PIB en 2020 al 2,4~\% en 2025, y la Comisión prevé el 2,4 y el 2,0~\% en 2026-2027 (cuadro~9). Con unos intereses del 2,4~\% del PIB, el saldo primario está equilibrado. La deuda, que pasó del 35,7~\% en 2007 a un máximo del 119,3~\% en 2020, está en el 100,5~\%. Pero el saldo total favorece la imagen: con la brecha en el +2,0~\%, el ciclo aporta en torno a 1,2 puntos, y el saldo ajustado por el ciclo es del $-3{,}6~\%$; sin medidas puntuales, el \textbf{saldo estructural} es del $-3{,}2~\%$ del PIB potencial, y el estructural primario, del $-0{,}8~\%$. Es decir, \uline{con la economía por encima de su potencial, las cuentas públicas aún tienen un déficit estructural del 3~\%}, por encima del 1,5~\% que fijan como referencia las reglas europeas. Los ingresos (42,9~\% del PIB) han subido casi 4 puntos desde 2019 (39,0~\%), y los gastos (45,3~\%) siguen 3 puntos por encima de los de antes de la pandemia (42,0~\%): parte del gasto extraordinario de 2020-2022 se ha quedado como gasto estructural.

La \textbf{inflación} ha ayudado a corto plazo, como predice la teoría: las bases tributarias son nominales y crecen con los precios (y, si los tramos del impuesto sobre la renta no se actualizan, la recaudación crece más que la renta), mientras los intereses solo suben a medida que se renueva la deuda (el tipo implícito pasa del 1,9 al 2,5~\% entre 2021 y 2025, pese a que el bono a 10 años llegó al 3,5~\% en 2023) y las pensiones y los salarios públicos se actualizan con retraso. A medio plazo el efecto se invierte: el gasto indexado y los intereses alcanzan a los ingresos.

\rotulo{b) Los rasgos de la política fiscal y su relación con el ciclo}
El tono se mide con la variación del saldo estructural primario, $\Delta b^{pe}$ (gráfico~13):
\begin{itemize}
\item en \textbf{2012-2013} la consolidación fue muy intensa (+2,9 y +2,5 puntos) con la brecha en su mínimo ($-8{,}7$ y $-9{,}6~\%$): una política \textbf{procíclica}, forzada por la pérdida de acceso a la financiación (diferencial de 435~pb);
\item en \textbf{2015-2019} fue expansiva, primero con la brecha negativa (contracíclica) y después con la brecha positiva ($-1{,}3$ puntos en 2019 con una brecha del +3,1~\%): procíclica de nuevo, porque no se aprovechó la expansión para reducir el déficit estructural y la economía llegó a la pandemia sin colchón;
\item en \textbf{2020-2022}, expansiva frente a la pandemia y a la crisis energética;
\item desde \textbf{2023}, contractiva con la brecha positiva (+1,3, +0,4 y +0,3 puntos en 2023-2025, y +0,4 en 2026-2027): por fin \textbf{contracíclica}, aunque \emph{tímida}, porque un déficit estructural del 3~\% que se reduce 3 o 4 décimas al año sigue siendo, por su nivel, una política holgada.
\end{itemize}
Los estabilizadores automáticos funcionan: en 2025 el ciclo mejora el saldo en 1,2 puntos, y en 2020 lo empeoró en 5.

\rotulo{c) La combinación con la política monetaria}
En 2026 la política monetaria pasa de algo expansiva a neutral para la zona del euro, pero sigue siendo expansiva para España (tipo real negativo, pregunta~7), y la fiscal es ligeramente contractiva. Las 2 van en la dirección correcta frente a una inflación alta y una brecha positiva: el \emph{policy mix} es \textbf{coherente}. Pero como la política monetaria es común y en España es más laxa que en la zona, \uline{la política fiscal es el único instrumento nacional para enfriar la demanda}, y debería apretar algo más de lo previsto; las ayudas frente a la energía, si las hay, deberían ser temporales y focalizadas en los hogares vulnerables. No siempre fue así: en 2018-2019 se combinó una política monetaria muy expansiva con una fiscal expansiva y una brecha positiva, y en 2022 las 2 seguían siendo expansivas con una inflación del 8~\%.

\rotulo{d) ¿Hay riesgos para la sostenibilidad de la deuda?}
La dinámica de la deuda es
\[
d_t-d_{t-1}=(g_t-\tau_t)+\frac{i_t-\gamma^N_t}{1+\gamma^N_t}\,d_{t-1}+\mathrm{ADD}_t .
\]
En 2025 el tipo implícito es del 2,5~\% y el crecimiento nominal, del 5,7~\% (gráfico~14): con $i-\gamma^N=-3{,}2$ puntos, el saldo primario que estabiliza la deuda es un \emph{déficit} del 3,1~\% del PIB. Con el saldo primario equilibrado la deuda debería haber bajado unos 3 puntos; bajó 0,9, lo que deja un ajuste entre déficit y deuda de unos 2 puntos (operaciones financieras que el cuadernillo no permite identificar). Entre 2020 y 2025 la deuda baja 18,8 puntos, y lo hace por el \textbf{efecto bola de nieve}, que la reduce en 31 puntos, frente a unos déficits primarios que suman 9: \uline{la deuda la ha bajado el crecimiento nominal, con la inflación incluida, no el ajuste fiscal}.

Por eso hay \textbf{riesgos apreciables a medio plazo}, aunque no inmediatos. El crecimiento nominal se moderará (5,1 y 4,5~\% previstos) y el tipo de la deuda nueva (el bono a 10 años está en el 3,9~\% en septiembre de 2026) ya supera al implícito (2,5~\%), de modo que el coste medio subirá a medida que se renueve la deuda. Si $i$ y $\gamma^N$ convergen hacia el 3,5~\%, harán falta superávits primarios para que la deuda baje. A ello se suman el gasto en pensiones y sanidad del envejecimiento (la tasa de dependencia sube del 24 al 31~\%), el de defensa y el fin de los fondos europeos. Lo mitigan que la deuda está en euros, un diferencial de solo 45~pb y una economía que crece.

\rotulo{e) Las reglas fiscales europeas}
Las reglas reformadas en 2024 mantienen las referencias del 3~\% de déficit y del 60~\% de deuda, pero se aplican con un solo indicador, una \textbf{senda de gasto primario neto} fijada en un plan fiscal-estructural de medio plazo (4 años, ampliables a 7 con reformas e inversiones), y con 2 salvaguardas: con una deuda superior al 90~\% del PIB, esta debe bajar al menos 1 punto al año de media, y el déficit estructural debe converger hacia un margen del 1,5~\%. Para España suponen 2 cosas: la deuda prevista (99,6 y 98,9~\%) baja 0,9 y 0,7 puntos al año, algo menos de lo que pide la salvaguarda, y el déficit estructural tiene que reducirse en torno a 1,7 puntos. Las reglas obligan, por tanto, a una \textbf{consolidación gradual y sostenida}, de algo más de medio punto al año, mejor por el lado del gasto corriente y ampliando bases que recortando inversión, y dejan margen para el gasto en defensa con la cláusula de escape nacional.
\end{examen}

\begin{complemento}{9.1 · Los saldos, del total al estructural primario}
\que{Qué se aplica.} Guía, V.3: las 5 variables del saldo, $b^{pe}=b-\varepsilon\cdot\mathrm{OG}-\text{puntuales}+\text{intereses}$ y $\text{esfuerzo}=\Delta b^{pe}$; «cuidado con el signo»; con déficits muy altos, mirar también el nivel \parencite{EFB2021FiscalStance}; los estabilizadores, $b-b^{ca}=\varepsilon\cdot\mathrm{OG}$.

\begin{tablacaja}
\centering\scriptsize\setlength{\tabcolsep}{3.5pt}
\caption{Del saldo total al esfuerzo fiscal (\% del PIB o del PIB potencial)}\label{tab:saldos}
\begin{tabular}{>{\raggedright\arraybackslash}p{3.5cm}*{9}{>{\centering\arraybackslash}p{0.93cm}}}
\toprule
\rowcolor{tceeazulpalido}
 & \textbf{2019} & \textbf{2020} & \textbf{2021} & \textbf{2022} & \textbf{2023} & \textbf{2024} & \textbf{2025} & \cellcolor{tceegrisclaro}\textbf{2026\,P} & \cellcolor{tceegrisclaro}\textbf{2027\,P} \\
\midrule
Saldo total, $b$ & $-3{,}1$ & $-9{,}9$ & $-6{,}7$ & $-4{,}6$ & $-3{,}3$ & $-3{,}2$ & $-2{,}4$ & $-2{,}4$ & $-2{,}0$ \\
Brecha, OG & 3,1 & $-8{,}5$ & $-3{,}3$ & 0,9 & 1,0 & 1,9 & 2,0 & 1,7 & 1,2 \\
Ajustado por el ciclo, $b^{ca}$ & $-4{,}9$ & $-4{,}8$ & $-4{,}7$ & $-5{,}1$ & $-3{,}9$ & $-4{,}3$ & $-3{,}6$ & $-3{,}5$ & $-2{,}7$ \\
\hspace*{0.8em}$\varepsilon$ implícita, $(b-b^{ca})/\mathrm{OG}$ & 0,59 & 0,60 & 0,61 & 0,58 & 0,64 & 0,61 & 0,59 & 0,60 & 0,60 \\
Estructural, $b^{e}$ & $-4{,}5$ & $-4{,}3$ & $-4{,}6$ & $-5{,}1$ & $-3{,}8$ & $-3{,}5$ & $-3{,}2$ & $-2{,}9$ & $-2{,}5$ \\
Estructural primario, $b^{pe}$ & $-2{,}2$ & $-2{,}1$ & $-2{,}5$ & $-2{,}8$ & $-1{,}5$ & $-1{,}0$ & $-0{,}8$ & $-0{,}4$ & 0,0 \\
\rowcolor{tceesalvia}
Esfuerzo, $\Delta b^{pe}$ & $-1{,}3$ & 0,1 & $-0{,}4$ & $-0{,}3$ & 1,3 & 0,4 & 0,3 & 0,4 & 0,4 \\
Tono y ciclo & \tiny exp. procícl. & \tiny neutral* & \tiny exp. contrac. & \tiny exp. procícl. & \tiny contr. contrac. & \tiny contr. contrac. & \tiny contr. contrac. & \tiny contr. contrac. & \tiny contr. contrac. \\
\bottomrule
\end{tabular}
\fuente{Cuadros~1 y~9 del cuadernillo (AMECO, primavera de 2026). Exp.: expansiva; contr.: contractiva; procícl./contrac.: procíclica o contracíclica según el signo de la brecha. (*) En 2020 las medidas de la pandemia se tratan como temporales y no entran en el saldo estructural (guía, V.3, nota): el saldo estructural no recoge el gran estímulo de ese año. Qué hay que ver: $\varepsilon\approx0{,}6$ en todos los años, la semielasticidad de la Comisión; la diferencia entre $b^{ca}$ y $b^{e}$ (0,4 puntos en 2025) son medidas puntuales que aumentan el déficit; $b^{pe}\approx b^e+\text{intereses}$ ($-3{,}2+2{,}4=-0{,}8$).}
\end{tablacaja}

\que{Error típico.} Juzgar el tono por el saldo total (que mejora 0,8 puntos en 2025 sobre todo por el ciclo y por los intereses) o leer al revés el signo: un $b^{pe}$ que pasa de $-1{,}0$ a $-0{,}8$ \emph{sube}: es contractivo (guía, V.3, subrayado).
\end{complemento}

\begin{complemento}{9.2 · La dinámica de la deuda, con números}
\que{Qué se aplica.} Guía, V.3 (ecuación de la dinámica de la deuda, saldo primario estabilizador, «la diferencia entre el tipo y el crecimiento nominal decide») y «La inflación y las cuentas públicas» (a corto plazo mejora la ratio; a medio plazo el efecto se invierte) \parencite{Blanchard2019Deuda}.

\que{Alternativa de redacción} de la (d) en 2 frases: «La deuda ha bajado 19 puntos desde 2020 sobre todo porque el PIB nominal ha crecido mucho más que el tipo de interés, no por el ajuste. Con el tipo marginal (3,9~\%) ya por encima del implícito (2,5~\%) y el crecimiento nominal moderándose, ese viento de cola se agota: sin superávits primarios la deuda dejará de bajar».

\begin{figuracaja}
\caption{España: variación de la deuda pública y sus componentes, 2021-2027}\label{fig:s08}
\grafico[0.72\linewidth]{s08-deuda}
\fuente{Cuadro~9 del cuadernillo; previsiones de 2026-2027 de la Comisión (tipo implícito previsto: intereses entre la deuda del año anterior). El ajuste déficit-deuda es el residuo. Qué hay que ver: la bola de nieve, muy negativa con la inflación de 2021-2023, pierde fuerza año a año; en 2026-2027 la deuda solo bajaría 0,7-0,9 puntos, menos que el punto anual de la salvaguarda.}
\end{figuracaja}

\begin{tablacaja}
\centering\scriptsize
\caption{Por qué bajó la deuda: descomposición de 2021-2025 (puntos del PIB)}\label{tab:deuda}
\begin{tabular}{>{\raggedright\arraybackslash}p{4.6cm}*{6}{>{\centering\arraybackslash}p{1.25cm}}}
\toprule
\rowcolor{tceeazulpalido}
 & \textbf{2021} & \textbf{2022} & \textbf{2023} & \textbf{2024} & \textbf{2025} & \textbf{2021-2025} \\
\midrule
Tipo implícito, $i$ (\%) & 1,9 & 2,2 & 2,4 & 2,5 & 2,5 & \\
Crecimiento nominal, $\gamma^N$ (\%) & 9,4 & 11,4 & 8,8 & 6,8 & 5,7 & \\
Déficit primario & 4,6 & 2,3 & 0,9 & 0,8 & 0,0 & 8,6 \\
Bola de nieve, $\frac{i-\gamma^N}{1+\gamma^N}d_{t-1}$ & $-8{,}1$ & $-9{,}5$ & $-6{,}4$ & $-4{,}3$ & $-3{,}1$ & $-31{,}5$ \\
Ajuste déficit-deuda (residuo) & $-0{,}1$ & 0,8 & 1,5 & $-0{,}4$ & 2,2 & 4,1 \\
\rowcolor{tceesalvia}
Variación de la deuda, $\Delta d$ & $-3{,}6$ & $-6{,}4$ & $-4{,}0$ & $-3{,}9$ & $-0{,}9$ & $-18{,}8$ \\
\bottomrule
\end{tabular}
\fuente{Cálculo propio con el cuadro~9 (\texttt{calculos.txt}). Ejemplo, 2025: $\frac{0{,}0249-0{,}0572}{1{,}0572}\times101{,}4=-3{,}1$; $\mathrm{ADD}=-0{,}9-0{,}0-(-3{,}1)=+2{,}2$.}
\end{tablacaja}

\begin{tablacaja}
\centering\scriptsize
\caption{Saldo primario que estabiliza una deuda del 100~\% del PIB según $i$ y $\gamma^N$}\label{tab:escenarios}
\begin{tabular}{>{\raggedright\arraybackslash}p{6.0cm}*{3}{>{\centering\arraybackslash}p{1.6cm}}}
\toprule
\rowcolor{tceeazulpalido}
\textbf{Escenario} & $\boldsymbol{i}$ & $\boldsymbol{\gamma^N}$ & \textbf{Saldo primario} \\
\midrule
2025 (observado) & 2,5 & 5,7 & $-3{,}0$ \\
2027 (crecimiento previsto, mismo tipo) & 2,5 & 4,5 & $-1{,}9$ \\
El tipo sube hacia el del bono; el crecimiento se modera & 3,5 & 4,0 & $-0{,}5$ \\
$i=\gamma^N$ & 3,5 & 3,5 & 0,0 \\
Escenario adverso & 4,0 & 3,0 & +1,0 \\
\bottomrule
\end{tabular}
\fuente{Cálculo propio: $\tau-g=\frac{i-\gamma^N}{1+\gamma^N}\,d$ con $d=100$ (guía, A.3.5). Qué hay que ver: el margen que da hoy $i<\gamma^N$ se va estrechando; para bajar la deuda 1 punto al año en el escenario central hará falta un superávit primario de algo más de medio punto.}
\end{tablacaja}

\end{complemento}

\begin{complemento}{9.3 · \emph{Policy mix} y reglas fiscales: cómo se razona}
\begin{figuracaja}
\caption{Los 4 cuadrantes del \emph{policy mix} (infografía de la guía)}\label{fig:policy-mix}
\grafico[0.46\linewidth]{policy-mix}
\fuente{Guía del dictamen, V.4. España en 2026: monetaria expansiva para España aunque neutral para la zona, fiscal ligeramente contractiva; con inflación y brecha positiva, el cuadrante deseable es el de «ambas contractivas».}
\end{figuracaja}

\que{Qué se aplica.} Guía, V.4: «el \emph{policy mix} es coherente si cada política contribuye a los objetivos de la otra, no si las 2 tienen el mismo signo»; en una unión monetaria «el país solo controla la fiscal, que tiene que compensar si su ciclo va por delante del de la zona». V.3: las reglas de 2024 \parencite{RegUE2024Preventivo,RegUE2024Correctivo} y la cláusula para el gasto en defensa \parencite{CE2025Defensa}; recomendaciones por desequilibrio (VI.2).

\que{Las salvaguardas, con números.} Deuda: 100,5 $\to$ 99,6 $\to$ 98,9, es decir, $-0{,}9$ y $-0{,}7$ puntos; la salvaguarda pide $-1{,}0$ de media en el periodo de ajuste. Déficit estructural: de $-3{,}2$ a $-1{,}5$ son 1,7 puntos; a 0,4 puntos al año (el esfuerzo previsto), unos 4 años. Como referencia externa, el FMI proponía un ajuste de en torno a 0,5 puntos del PIB al año entre 2026 y 2030 \parencite{FMI2026ArticuloIV}.

\que{Errores típicos.} Decir que España «incumple» las reglas por tener una deuda del 100~\%: el valor de referencia del 60~\% no se exige de golpe, sino una senda. Y recomendar una consolidación brusca: con el choque energético frenando la demanda, mejor gradual y creíble, anunciada pronto (guía, V.3).

\que{Qué añadir con más tiempo.} Las necesidades brutas de financiación, la vida media de la deuda y quién la tiene (no vienen en el cuadernillo; guía V.3); la calidad del ajuste (proteger la inversión pública y la educación; ampliar bases antes que subir tipos). \que{Qué recortar en el examen.} La lista de etapas del tono (basta: procíclica en 2012-2013 y 2019, contracíclica y tímida desde 2023) y la historia del \emph{policy mix} anterior a 2026.
\end{complemento}

% Pregunta 10. Desequilibrios, riesgos y recomendaciones (bloque VI) · peso alto
\section*{Pregunta 10. Desequilibrios, riesgos y recomendaciones}

\begin{enunciado}
\textbf{10.} Identifique los principales \textbf{desequilibrios y vulnerabilidades} de la economía española y los riesgos al alza y a la baja que rodean las previsiones para 2026 y 2027. Teniendo en cuenta las tendencias demográficas, ¿qué \textbf{reformas estructurales y políticas de oferta} recomendaría para sostener el crecimiento a medio y largo plazo?

\smallskip
\textbf{Peso alto} · \tiempo{21} · guía: VI.1 (fortalezas, desequilibrios y medidas; naturaleza, origen y gravedad), VI.2 (recomendaciones por desequilibrio), V.5 (políticas de oferta), D.8 (matriz), B.1 (demografía) · todo el cuadernillo.
\end{enunciado}

\begin{examen}
\textit{España afronta el final de la década con fortalezas que no tenía en 2007 (superávit exterior, sector privado desapalancado, banca rentable, empleo en máximos) y con 3 desequilibrios de fondo, por este orden: una productividad estancada, unas cuentas públicas con déficit estructural y deuda alta, y un paro estructural elevado; a ellos se suma, en 2026, un choque de precios que puede volverse persistente. El envejecimiento hace más urgente corregirlos, porque el crecimiento no podrá apoyarse indefinidamente en sumar trabajadores.}

\rotulo{a) Desequilibrios y vulnerabilidades}
Las \textbf{fortalezas} son claras: un crecimiento superior al de la zona del euro, la tasa de empleo más alta de la serie, un superávit exterior que se mantiene desde 2012, la deuda privada en su mínimo (hogares, 42,9~\%; empresas, 62,0~\% del PIB), una banca rentable y una prima de riesgo de 45~pb. Los desequilibrios, ordenados por su importancia:

\textbf{1.}~\textbf{Productividad estancada}: el PIB por hora crece un 0,33~\% anual desde 2019, el gasto en I+D es dos tercios del europeo y la inversión en maquinaria apenas crece. Es estructural y persistente, y es el más importante porque \uline{solo la productividad eleva la renta por habitante}, que desde 2007 crece menos que en la zona del euro.


\textbf{2.}~\textbf{Finanzas públicas}: un déficit estructural del 3,2~\% con la brecha positiva y una deuda del 100,5~\%, que ha bajado gracias al crecimiento nominal. Es un problema de \emph{stock}, persistente, que se agravará con el envejecimiento y con la normalización del coste de la deuda.


\textbf{3.}~\textbf{Paro estructural}: el 10,5~\%, el doble que en la zona del euro, con un paro juvenil del 24,9~\% y muy pocas vacantes: un problema de emparejamiento y de cualificación más que de demanda.


\textbf{4.}~\textbf{Inflación y costes}: coyuntural, pero con riesgo de enquistarse; la inflación del 5,0~\% y unos costes laborales que crecen más que en la zona erosionan la competitividad ganada.

Los 2 primeros son problemas de \emph{stock} y de oferta, que no se corrigen solos con el ciclo; el tercero es en parte la histéresis de las crisis pasadas; el cuarto es de flujo, pero los problemas de flujo que persisten acaban siendo de \emph{stock}, y por eso conviene atajarlo pronto. Y 3 \textbf{vulnerabilidades}: la dependencia energética (el déficit de bienes llegó al 4,4~\% del PIB en 2022 y el superávit corriente se reducirá en 2026), la vivienda (precios a 2 dígitos con una oferta que no responde) y una posición exterior todavía por debajo del umbral europeo, aunque cada vez menos preocupante.

\rotulo{b) Riesgos de las previsiones}
Las previsiones de la Comisión (PIB del 2,4 y el 1,9~\%; inflación del 3,0~\% en 2026) son de primavera, anteriores a buena parte del choque. Los riesgos para el \textbf{crecimiento} son, en 2026, equilibrados: el ritmo del primer semestre (0,6-0,7~\% trimestral) permitiría superar la previsión, pero el petróleo, por encima de los 110 dólares en septiembre, resta renta real a los hogares y márgenes a las empresas, y empeora la balanza de bienes (en el choque de 2022 el déficit de bienes pasó del 1,7 al 4,4~\% del PIB). En 2027 dominan los riesgos \textbf{a la baja}: el encarecimiento del crédito (el Euríbor ha subido 1 punto y gran parte de las hipotecas son a tipo variable), la debilidad de los socios por los aranceles, la consolidación fiscal pendiente, el final de los fondos europeos y un posible ajuste de la vivienda. \textbf{Al alza}: la llegada de población, el turismo, un ahorro de los hogares todavía alto (11,9~\%) que da margen al consumo, y la inversión en intangibles. Para la \textbf{inflación}, los riesgos son claramente al alza: la media de enero a septiembre (3,6~\%) ya supera la previsión anual, y la subyacente no cede.

\rotulo{c) Reformas estructurales y políticas de oferta}
\textbf{A corto plazo}, la prioridad es que el choque energético no se convierta en inflación persistente: una política fiscal que no añada demanda (ayudas, si las hay, temporales y focalizadas en los hogares vulnerables y que mantengan el incentivo a ahorrar energía), un pacto de rentas y una política macroprudencial preparada para frenar el crédito hipotecario si se acelera. La política monetaria es común y no puede ajustarse a España.

\textbf{A medio y largo plazo}, el punto de partida es la \textbf{demografía}: la tasa de dependencia de los mayores sube del 24,0 al 31,2~\% entre 2007 y 2025, y el peso de la población en edad de trabajar ya no suma al crecimiento (cuadro~2~bis), que se ha apoyado en la inmigración. Para sostenerlo hace falta que crezca la productividad y que trabaje más gente. Las recomendaciones, de la más a la menos importante, con su objetivo y su riesgo:
\begin{itemize}
\item \textbf{Productividad}: capital humano (menos abandono escolar, formación profesional dual, recualificación), I+D hasta acercarse a la media europea (incentivos fiscales, colaboración entre universidades y empresas, capital riesgo) y un entorno que permita a las empresas crecer (menos trabas regulatorias ligadas al tamaño, unidad de mercado, competencia en los servicios). Objetivo: un crecimiento intensivo. Riesgo: sus efectos tardan años.
\item \textbf{Mercado de trabajo}: políticas activas de empleo focalizadas en jóvenes y parados de larga duración y evaluadas, mejor intermediación de los servicios públicos de empleo, consolidar la menor temporalidad y salarios ligados a la productividad, con una negociación colectiva que tenga en cuenta la situación de cada empresa. Objetivo: bajar el paro estructural hacia el europeo. Riesgo: la eficacia de las políticas activas depende de su diseño y su evaluación.
\item \textbf{Demografía y pensiones}: integrar bien la inmigración (reconocimiento de cualificaciones, formación), elevar la actividad de los mayores y de las mujeres (jubilación flexible, servicios de cuidados) y asegurar la sostenibilidad de las pensiones ligándolas a la evolución demográfica. Objetivo: que la población activa siga creciendo cuando se jubilen las generaciones más numerosas.
\item \textbf{Consolidación fiscal gradual}: algo más de medio punto del PIB al año, con revisiones del gasto, menos beneficios fiscales y bases más amplias, protegiendo la inversión pública; el plan fiscal-estructural le da credibilidad. Riesgo: frenar la demanda si se hace deprisa.
\item \textbf{Vivienda}: actuar sobre la \emph{oferta} (suelo, licencias, alquiler asequible) más que sobre la demanda, porque la vivienda cara empeora el emparejamiento laboral y la movilidad.
\item \textbf{Energía}: diversificar las fuentes, acelerar las renovables, reforzar las interconexiones y la eficiencia, para reducir la exposición a choques como los de 2022 y 2026.
\end{itemize}
En conclusión, España tiene hoy el margen que no tuvo en 2008 (sin desequilibrios privados ni exteriores) para \uline{aprovechar la expansión y corregir los desequilibrios públicos y de productividad}; cómo lo haga decidirá si su crecimiento se traduce en renta por habitante.
\end{examen}

\begin{complemento}{10.1 · La matriz de fortalezas y debilidades, rellena}
\que{Qué se aplica.} Guía, D.8: la matriz se va rellenando durante el examen en una hoja aparte y sirve para 3 cosas: anticipar el diagnóstico en la pregunta~1, no contradecirse y contestar la última pregunta. VI.1: \emph{qué va bien, qué hay que cambiar (2 o 3 desequilibrios ordenados) y cómo}; las 3 preguntas del FMI (naturaleza, origen y gravedad) \parencite{FMI2013CaseStudy}.

\begin{tablacaja}
\centering\scriptsize
\renewcommand{\arraystretch}{0.98}\setlength{\tabcolsep}{4pt}
\caption{Matriz de fortalezas y debilidades de España (examen de ejemplo)}\label{tab:matriz-espana}
\begin{tabular}{>{\raggedright\arraybackslash}p{2.0cm}>{\raggedright\arraybackslash}p{3.3cm}>{\raggedright\arraybackslash}p{3.6cm}>{\centering\arraybackslash}p{0.7cm}>{\centering\arraybackslash}p{0.8cm}>{\raggedright\arraybackslash}p{3.4cm}}
\toprule
\rowcolor{tceeazulpalido}
\textbf{Bloque} & \textbf{\textcolor{tceeverde}{Fortaleza}} & \textbf{\textcolor{tceerojo}{Debilidad}} & \textbf{Peso} & \textbf{Tend.} & \textbf{Recomendación} \\
\midrule
Perfil cíclico & PIB +10,0~\% en 2019-2025 (zona: +6,8~\%) & Crecimiento extensivo: PIB por habitante +5,0~\% & ** & $\nearrow$ & Políticas de oferta (productividad) \\
Trabajo & Empleo 67,0~\%; temporalidad 15,4~\% & Paro 10,5~\%; juvenil 24,9~\%; PIB por hora +0,33~\%/año & *** & $\nearrow$ & Políticas activas, emparejamiento, capital humano \\
Capital y crédito & Deuda privada en mínimos; intangibles +47~\% & Maquinaria +2,7~\%; vivienda +12,7~\% & * & $\rightarrow$ & Oferta de vivienda; incentivos a invertir \\
Precios & Salida rápida de 2022-2023 & IPCA 5,0~\% (zona 3,8); CLU +4,4~\% & ** & $\searrow$ & Pacto de rentas; fiscal que no empuje \\
No precio & Cuota mundial estable; servicios & I+D 1,50~\% (UE 2,26) & * & $\rightarrow$ & I+D, tamaño, diversificar mercados \\
Exterior & Superávit desde 2012; NENDI $-9{,}7~\%$ & PIIN $-42{,}7~\%$; energía & * & $\nearrow$ & Política energética \\
Monetaria & Transmisión rápida; diferencial 45~pb & Tipo real más bajo en España & * & $\rightarrow$ & (común; compensa la fiscal) \\
Banca & ROE 14,2~\%; dudosos 2,4~\% & CET1 13,9~\% (UE 16,4) & * & $\nearrow$ & Colchones; límites LTV/LTI si hace falta \\
Fiscal & Saldo primario equilibrado & Déficit estructural 3,2~\%; deuda 100,5~\% & *** & $\nearrow$ & Consolidación gradual con la senda de gasto \\
Demografía & Inmigración (extranjeros 14,1~\%) & Dependencia 31,2~\% & ** & $\searrow$ & Integración, actividad, pensiones \\
\bottomrule
\end{tabular}
\fuente{Elaboración propia con el cuadernillo; plantilla de la guía (D.8). Peso: importancia del desequilibrio (*** máxima). Tendencia: $\nearrow$ mejora, $\rightarrow$ estable, $\searrow$ empeora.}
\end{tablacaja}
\end{complemento}

\begin{complemento}{10.2 · Los umbrales oficiales: qué se enciende en España}

\begin{tablacaja}
\centering\scriptsize
\caption{España frente a los umbrales del procedimiento de desequilibrios y del Tratado}\label{tab:pdm}
\begin{tabular}{>{\raggedright\arraybackslash}p{4.4cm}>{\centering\arraybackslash}p{2.6cm}>{\centering\arraybackslash}p{3.2cm}>{\centering\arraybackslash}p{2.4cm}}
\toprule
\rowcolor{tceeazulpalido}
\textbf{Indicador} & \textbf{Umbral} & \textbf{España} & \textbf{¿Alerta?} \\
\midrule
Saldo corriente, media de 3 años & $-4$ / $+6~\%$ & $(2{,}7+3{,}3+3{,}2)/3=3{,}1$ & \textcolor{tceeverde}{no} \\
PIIN & $-35~\%$ & $-42{,}7~\%$ & \textcolor{tceerojo}{sí} (mejorando) \\
Tipo de cambio efectivo real, 3 años & $\pm3~\%$ & $-0{,}9~\%$ (2022-2025) & \textcolor{tceeverde}{no} \\
CLU nominales, 3 años & $+9~\%$ & $\approx+13{,}9~\%$ (2023-2025)* & \textcolor{tceerojo}{sí} \\
Precio de la vivienda (nominal) & $+9~\%$ & $+12{,}7~\%$ & \textcolor{tceerojo}{sí} \\
Deuda de hogares / de empresas & 55 / 85~\% & 42,9 / 62,0~\% & \textcolor{tceeverde}{no} \\
Tasa de paro & 10~\% & 10,5~\% & \textcolor{tceerojo}{sí} (por poco) \\
Saldo público (Tratado) & $-3~\%$ & $-2{,}4~\%$ & \textcolor{tceeverde}{no} \\
Deuda pública (Tratado) & 60~\% & 100,5~\% & \textcolor{tceerojo}{sí} \\
\bottomrule
\end{tabular}
\fuente{Umbrales: \textcite{CE2025AMR2026,TFUE2012Protocolo12} (guía, A.1); datos: cuadros 3, 4, 6, 7 y 9. (*) El umbral se refiere a los CLU por hora; con los CLU por persona del cuadernillo, $1{,}059\times1{,}030\times1{,}044-1\approx13{,}9~\%$, una aproximación. Qué hay que ver: las alertas están en los costes, la vivienda, el paro y la deuda pública: exactamente los desequilibrios de la respuesta. Son señales, no obligaciones (guía, A).}
\end{tablacaja}
\begin{figuracaja}
\caption{Los umbrales oficiales de un vistazo (infografía de la guía)}\label{fig:umbrales}
\grafico[0.6\linewidth]{esquema-umbrales-pdm}
\fuente{Guía del dictamen, A (umbrales del procedimiento de desequilibrios y del Tratado). La tabla anterior sitúa a España en cada uno.}
\end{figuracaja}
\end{complemento}

\begin{complemento}{10.3 · Riesgos, recomendaciones que responden a los datos y qué recortar}
\que{Del dato a la recomendación} (guía, VI.1: «puntúan las recomendaciones que responden a los datos del cuadro»): cada medida de la respuesta sale de un indicio (gráfico~\ref{fig:dato-recomendacion}). Ejemplo: \emph{PIB por hora +0,33~\% e I+D del 1,50~\%} $\to$ productividad estancada $\to$ I+D y tamaño empresarial $\to$ objetivo: crecimiento intensivo $\to$ riesgo: tarda.

\begin{figuracaja}
\caption{Del dato a la recomendación (infografía de la guía)}\label{fig:dato-recomendacion}
\grafico[0.56\linewidth]{esquema-dato-recomendacion}
\fuente{Guía del dictamen, VI.1. Qué hay que ver: siempre el mismo camino, del indicio al desequilibrio, a la medida, a su objetivo y a su riesgo.}
\end{figuracaja}

\begin{tablacaja}
\centering\scriptsize
\caption{Riesgos de las previsiones de 2026-2027 y su indicio en el cuadernillo}\label{tab:riesgos}
\begin{tabular}{>{\raggedright\arraybackslash}p{1.6cm}>{\raggedright\arraybackslash}p{5.2cm}>{\raggedright\arraybackslash}p{7.0cm}}
\toprule
\rowcolor{tceeazulpalido}
\textbf{Sentido} & \textbf{Riesgo} & \textbf{Indicio} \\
\midrule
A la baja & Petróleo caro y persistente & Brent de 62,5 a 114,2 dólares (G7); bienes $-4{,}4~\%$ del PIB en el choque de 2022 (C7) \\
 & Endurecimiento financiero & Depósito 2,50~\%; Euríbor +98~pb desde diciembre de 2025; hipotecas 3,04~\% (G12) \\
 & Aranceles y socios & Exportaciones a EE.~UU. $-11~\%$ (G9) \\
 & Fin de los fondos y consolidación & Cuenta de capital $\approx1~\%$ del PIB (C7); esfuerzo +0,4 en 2026-2027 (C9) \\
 & Vivienda & Precio +12,7~\% con poca inversión residencial (C1, C4) \\
Al alza & Población y turismo & Población +1,0~\% al año; viajes 4,3~\% del PIB (C3, C7) \\
 & Consumo & Ahorro de los hogares 11,9~\%, el doble que en 2007 (C4) \\
 & Ritmo de 2026 & 0,6-0,7~\% trimestral en el primer semestre (G1): $\approx$2,5~\% anual \\
Inflación (al alza) & Segunda ronda & Media de 2026 ya 3,6~\% frente a 3,0 previsto; subyacente 3,5~\% (G6) \\
\bottomrule
\end{tabular}
\fuente{Cuadernillo (C: cuadro; G: gráfico).}
\end{tablacaja}

\que{Errores típicos} (guía, VI.1 y D.6): recomendaciones genéricas que valen para cualquier país («hacer reformas estructurales»); recetas tajantes o demasiado precisas; olvidar la demografía, que el enunciado pide expresamente; no ordenar los desequilibrios. \que{Alternativa de redacción} del cierre: «España tiene que aprovechar una expansión sin desequilibrios privados para corregir los públicos y elevar la productividad, antes de que el envejecimiento reduzca el crecimiento de la población activa». \que{Qué añadir con más tiempo.} El caso de España del FMI (guía, C.2): ajuste fiscal de unos 0,5 puntos al año entre 2026 y 2030, políticas activas y de vivienda, y reformas de productividad \parencite{FMI2026ArticuloIV}. \que{Qué recortar en el examen.} Las vulnerabilidades (1 frase) y 2 de las 6 recomendaciones (energía y vivienda, ya tratadas en las preguntas 3, 4 y 8).
\end{complemento}


\bibliografia
\end{document}
